% Generated mechanically from frozen input p01/ko/fields.tex.
% Do not edit this generated file; edit the bound locale source instead.

% OK, start here.
%

\setcounter{chapter}{8}
\chapter{체}



\phantomsection
\label{fields-section-phantom}



\section{서론}
\label{fields-section-introduction}

\noindent
이 장에서는 체 이론을 다룬다. {\it 체}란 모든 영이 아닌 원소가
가역인 영환이 아닌 환임을 상기하자. 동치로, 영이 아닌 모든 원소가
단원이므로 체의 이데알은 $(0)$과 $(1)$뿐이다. 따라서 뒤에서 전개할
이론의 많은 부분에서 체는 가장 단순한 경우가 된다.

\medskip\noindent
예를 들어 체의 사상은 모두 단사이므로, 체 확대 이론은 표준적인
가환대수와는 성격이 다르다. 그럼에도 환에 관한 문제를 체에 관한
문제로 환원할 수 있는 경우가 많다. 예를 들어 모든 정역은 하나의 체
(그 분수체)에 매입될 수 있고, 모든 {\it 국소환}(즉 유일한 극대
이데알을 갖는 환; 이 용어는 아직 정의하지 않았다)에는 그 잉여체
(즉 극대 이데알에 의한 몫)이 결부된다. 그러므로 체 확대에 관한
지식은 유용할 것이다.




\section{기본 정의}
\label{fields-section-definitions}

\noindent
이 장을 가환대수를 논하는 장보다 앞에 두었으므로, 뒤의 대수 장에서
더 자세히 다룰 몇 가지 기본 정의를 여기에서 먼저 도입해야 한다.

\begin{definition}
\label{fields-definition-field}
{\it 체}란 모든 영이 아닌 원소가 가역인 영환이 아닌 환이다.
주어진 체의 {\it 부분체}란 그 자체가 체인 부분환이다.
\end{definition}

\noindent
체 $k$에 대하여 부분집합 $k \setminus \{0\}$을 $k^*$로 쓴다.
이는 환 $R$의 가역원들의 군을 나타내는 통상적인 기호 $R^*$를
일반화한 것이다.

\begin{definition}
\label{fields-definition-domain}
{\it 정역} 또는 {\it 정수환}이란 영인자가 $0$뿐인 영환이 아닌
환이다.
\end{definition}



\section{체의 예}
\label{fields-section-examples}

\noindent
먼저 체의 몇 가지 예를 들자. 환 $R$과 이데알 $I \subset R$에
대하여 $R/I$가 체일 필요충분조건은 $I$가 극대 이데알인 것임을
상기하자.

\begin{example}[유리수]
\label{fields-example-rational-numbers}
유리수들은 하나의 체를 이룬다. 이를 {\it 유리수체}라 하며
$\mathbf{Q}$로 나타낸다.
\end{example}

\begin{example}[소체]
\label{fields-example-prime-field}
$p$가 소수이면 $\mathbf{Z}/(p)$는 체이고, 이를 $\mathbf{F}_p$로
나타낸다. 실제로 $(p)$는 $\mathbf{Z}$의 극대 이데알이다. 따라서
체는 유한할 수도 있으며, $\mathbf{F}_p$는 $p$개의 원소를 갖는다.
\end{example}

\begin{example}
% P01-E0235: frozen stable label misspelling retained.
\label{fields-example-quotient-polymial-ring}
주 아이디얼 정역에서 기약원 하나가 생성하는 이데알은 극대이다.
이제 $k$가 체이면 다항식환 $k[x]$는 주 아이디얼 정역이다. 따라서
$P \in k[x]$가 기약다항식(즉 차수가 더 작은 항들의 곱으로 인수분해되지
않는 비상수 다항식)이면 $k[x]/(P)$는 체이다. 이 체는 자연스럽게
$k$의 한 복사본을 포함한다. 이는 체를 구성하는 매우 일반적인
방법이다. 예를 들어 복소수체 $\mathbf{C}$는
$\mathbf{R}[x]/(x^2 + 1)$로 구성할 수 있다.
\end{example}

\begin{example}[분수체]
\label{fields-example-quotient-field}
정역 $A$가 주어지면, $\mathbf{Q}$를 $\mathbf{Z}$로부터 구성하는
것과 정확히 같은 방식으로 $A$에서 구성한 체 $F$로의 매입
$A \to F$가 있음을 상기하자. 형식적으로 $F$의 원소들은 분수
$a/b$의 동치류이며, 여기서 $a, b \in A$, $b \not = 0$이다.
통상대로 $a/b = a'/b'$일 필요충분조건은 $ab' = ba'$이다.
체 $F$를 $A$의 {\it 몫체}, {\it 분수들의 체}, 또는 {\it 분수체}라
한다. 분수체에는 다음 보편 성질이 있다. 체 $K$로 가는 단사 환 사상
$\varphi : A \to K$가 주어지면 다음 도표를 가환하게 하는 유일한
사상 $\psi : F \to K$가 존재한다.
$$
\xymatrix{
F \ar[r]_\psi & K \\
A \ar[u] \ar[ru]_\varphi
}
$$
실제로 그러한 사상을 정의하는 법은 명백하다. $\psi(a/b) =
\varphi(a)\varphi(b)^{-1}$로 두면 된다. $b \not = 0$일 때
$\varphi(b) \not = 0$임은 $\varphi$의 단사성으로 보장된다.
\end{example}

\begin{example}[유리함수체]
\label{fields-example-field-of-rational-functions}
$k$가 체이면 $k$ 위의 유리함수체 $k(x)$를 생각할 수 있다. 이는
다항식환 $k[x]$의 분수체이다. 다시 말해, 자명한 동치관계를 준
$F, G \in k[x]$, $G \not = 0$인 몫 $F/G$들의 집합이다.
\end{example}

\begin{example}
\label{fields-example-field-of-meromorphic-functions}
$X$를 리만 곡면이라 하고, $X$ 위의 유리형 함수들의 집합을
$\mathbf{C}(X)$로 나타내자. 함수의 곱셈과 덧셈에 관하여
$\mathbf{C}(X)$는 환이다. 실제로 $\mathbf{C}(X)$는 체임을 알 수
있다. 곧, 영이 아닌 함수 $f(z)$가 유리형이면 $1/f(z)$도 유리형이다.
예를 들어 $S^2$를 리만 구면이라 하자. 복소해석학에 의해 유리형
함수환 $\mathbf{C}(S^2)$가 유리함수체 $\mathbf{C}(z)$임을 안다.
\end{example}



\section{벡터공간}
\label{fields-section-vector-spaces}

\noindent
체가 다루기 좋은 한 가지 이유는 체 위 가군, 즉 벡터공간의 이론이
매우 단순하기 때문이다.

\begin{lemma}
\label{fields-lemma-vector-space-is-free}
$k$가 체이면 모든 $k$-가군은 자유이다.
\end{lemma}

\begin{proof}
실제로 선형대수학에 의해 $k$-가군, 즉 벡터공간 $V$에는
{\it 기저} $\mathcal{B} \subset V$가 있고, 이 기저는
$\mathcal{B}$ 위의 자유 벡터공간에서 $V$로 가는 동형사상을 정한다.
\end{proof}

\begin{lemma}
\label{fields-lemma-field-semi-simple}
체 위 가군들의 모든 완전열은 분할된다.
\end{lemma}

\begin{proof}
모든 벡터공간이 사영 가군이므로 이는
Lemma \ref{fields-lemma-vector-space-is-free}에서 따른다.
\end{proof}

\noindent
뒤의 장에서 전개할 이론이 체에 관해 많은 말을 하지 않는 또 다른
이유도 가군이 이처럼 단순하게 행동하기 때문이다. Lemma
\ref{fields-lemma-field-semi-simple}는 $k$가 체일 때 $k$-가군들의
{\it 범주}에 관한 명제임에 유의하자. 완전성이라는 개념은 본질적으로
화살표론적, 즉 순전히 범주론적인 개념을 사용하며, 실제로 이른바
{\it 아벨 범주} 안에서 서술할 수 있기 때문이다.

\medskip\noindent
이제부터 체 위 가군의 연구는 선형대수학이고 체의 이데알 이론은
그다지 흥미롭지 않으므로, 이 장의 실제 주제인 체의 {\it 확대}를
연구한다.


\section{체의 표수}
\label{fields-section-more-fields}

\noindent
환의 범주에는 {\it 초기 대상} $\mathbf{Z}$가 있다. 즉 모든 환
$R$에는 $\mathbf{Z}$에서 오는 사상이 정확히 하나 있다. 체의
범주에는 그러한 초기 대상이 없다. 그럼에도 다음과 같은 대상들의
족이 있다. 모든 체는 그들 중 정확히 하나에 의해 정확히 한 가지
방식으로 사상되고, 나머지 대상들에 의해서는 어떤 방식으로도
사상되지 않는다.
% P01-E0236: frozen malformed source-target wording recorded in the sidecar.

\medskip\noindent
$F$를 체라 하자. $F$를 환으로 보아 환 사상
$f : \mathbf{Z} \to F$를 얻는다. 이 환 사상의 상은 체의 부분환이므로
정역이고, 따라서 $f$의 핵은 $\mathbf{Z}$의 소 이데알이다. 그러므로
$f$의 핵은 $(0)$이거나 어떤 소수 $p$에 대한 $(p)$이다.

\medskip\noindent
첫 번째 경우에는 $f$가 단사이므로 $\mathbf{Z}$를 $F$의 부분환으로
본다. 더구나 $F$의 영이 아닌 모든 원소는 가역이므로
$p, q \in \mathbf{Z}$, $q \not = 0$에 대하여 $p/q \in F$라고 말할
수 있다. 따라서 이 경우 $\mathbf{Q}$를 $F$의 부분환으로 볼 수 있고
실제로 그렇게 한다. 이것이 이 경우 $F$의 가장 작은 부분체임은 쉽게
알 수 있다.

\medskip\noindent
두 번째 경우, 즉 $\Ker(f) = (p)$이면
$\mathbf{Z}/(p) = \mathbf{F}_p$가 $F$의 부분환이다. 명백히 이는
$F$의 가장 작은 부분체이다.

\medskip\noindent
이와 같이 논하면 모든 체에는 가장 작은 부분체가 있으며, 이는
$\mathbf{Q}$이거나 어떤 소수 $p$에 대해 유한하고
$\mathbf{F}_p$와 같다.
% P01-E0237: frozen “finite equal to” wording recorded in the sidecar.

\begin{definition}
\label{fields-definition-characteristic}
체 $F$의 {\it 표수}는 $\mathbf{Z} \subset F$이면 $0$이고,
$F$에서 $p = 0$인 소수 $p$가 있으면 그 소수이다. $F$의
{\it 소부분체}란 $F$의 가장 작은 부분체로, 표수가 영이면
$\mathbf{Q} \subset F$이고 표수가 $p > 0$이면
$\mathbf{F}_p \subset F$이다.
\end{definition}

\noindent
$E \subset F$가 부분체이면 $E$의 표수와 $F$의 표수가 같음은
쉽게 알 수 있다.

\begin{example}
\label{fields-example-characteristic}
$\mathbf{F}_p$의 표수는 $p$이고 $\mathbf{Q}$의 표수는 $0$이다.
\end{example}


\section{체 확대}
\label{fields-section-extensions}

\noindent
일반적으로 우리는 체 자체보다 체의 {\it 확대}에 더 관심이 있다.
이는 고정된 환 위에서 환 자체가 아니라 {\it 대수}를 연구하는 것과
비슷할 수 있다. 다음 결과에 의해 ``다른 체 위의 체''라는 개념이
바로 체 확대라는 개념으로 귀착된다는 점이 체에서는 편리하다.

\begin{lemma}
\label{fields-lemma-field-maps-injective}
$F$가 체이고 $R$이 영환이 아닌 환이면, 모든 환 준동형
$\varphi : F \to R$은 단사이다.
\end{lemma}

\begin{proof}
실제로 $a \in \Ker(\varphi)$를 영이 아닌 원소라 하자. 그러면
$\varphi(1) = \varphi(a^{-1} a) = \varphi(a^{-1}) \varphi(a) = 0$이다.
따라서 $1 = \varphi(1) = 0$이고 $R$은 영환이다.
\end{proof}

\begin{definition}
\label{fields-definition-extension}
체 $F$가 체 $E$에 포함되어 있으면 $E$를 $F$의 {\it 체 확대}라고
한다. $E$가 $F$의 확대임을 $E/F$로 나타낸다.
\end{definition}

\noindent
따라서 $F, F'$가 체이고 $F \to F'$가 임의의 환 준동형이면,
Lemma \ref{fields-lemma-field-maps-injective}에 의해 이는 단사이고, 약간
용어를 남용하여 $F'$를 $F$의 확대라고 볼 수 있다. 달리 말해
$F$의 체 확대란 우연히 체인 $F$-대수일 뿐이다. 일반적인 환의
경우에는 환 준동형이 반드시 단사인 것은 아니므로 상황이 전혀 다르다.

\medskip\noindent
$k$를 체라 하자. $k$의 체 확대들의 {\it 범주}가 있다. 이 범주의
대상은 확대 $E/k$, 즉 (반드시 단사인) 체의 사상
$$
k \to E,
$$
이고, 확대 $E/k$와 $E'/k$ 사이의 사상은 $k$-대수 사상
$E \to E'$이다. 달리 말해 이는 가환 도표
$$
\xymatrix{
E \ar[rr] & & E' \\
& k \ar[ru] \ar[lu] &
}
$$
이다. $k$의 확대들의 범주에서 $E \to E'$로 가는 사상들의 집합은
$\Mor_k(E, E')$로 나타낸다.
% P01-E0238: frozen malformed source-target wording recorded in the sidecar.

\begin{definition}
\label{fields-definition-tower}
체의 {\it 탑} $E_n/E_{n - 1}/\ldots/E_0$은 체 확대들의 열
$E_n/E_{n - 1}$, $E_{n - 1}/E_{n - 2}$, $\ldots$, $E_1/E_0$로
이루어진다.
\end{definition}

\noindent
체 확대의 몇 가지 예를 들자.

\begin{example}
\label{fields-example-monogenic-extension}
$k$를 체라 하고 $P \in k[x]$를 기약다항식이라 하자.
$k[x]/(P)$가 체임을 이미 보았다
(Example \ref{fields-example-quotient-polymial-ring}). 이는 자명한 방식으로
$k$-대수이기도 하므로 $k$의 확대이다.
\end{example}

\begin{example}
\label{fields-example-field-of-meromorphic-functions-extension-C}
$X$가 리만 곡면이면 유리형 함수체 $\mathbf{C}(X)$
(Example \ref{fields-example-field-of-meromorphic-functions})는
$\mathbf{C}$의 확대체이다. 실제로 $\mathbf{C}$의 각 원소는 $X$ 위의
유리형이고, 더 나아가 정칙인 상수함수를 유도한다.
\end{example}

\noindent
$F/k$를 체 확대라 하고 $S \subset F$를 임의의 부분집합이라 하자.
그러면 $S$를 포함하는 $F$의 {\it 가장 작은} 부분확대, 즉 $k$를
포함하는 $F$의 부분체가 있다. 이를 보려면 $S$와 $k$를 포함하는
$F $의 부분체들의 족을 생각하여 그 교집합을 취하라. 이것이 체임을
확인할 수 있다. 실제로 표준적인 논증에 의해 이것은 $k$와 $S$의
원소들을 사용한 유한 번의 기본 대수 연산(덧셈, 곱셈, 뺄셈, 나눗셈)으로
얻을 수 있는 $F$의 원소들의 집합이다.

\begin{definition}
\label{fields-definition-generated-by}
$k$를 체라 하자. $F/k$가 체 확대이고 $S \subset F$이면, $k$와
$S$를 포함하는 $F$의 가장 작은 부분체를 $k(S)$로 쓴다. $S$가
체 확대 $k(S)/k$를 {\it 생성한다}고 한다. $S = \{\alpha\}$가
한원소집합이면 $k(\{\alpha\})$ 대신 $k(\alpha)$로 쓴다.
$F = k(S)$인 유한 부분집합 $S \subset F$가 존재하면 $F/k$를
{\it 유한생성 체 확대}라고 한다.
\end{definition}

\noindent
예를 들어 $\mathbf{C}$는 $\mathbf{R}$ 위에서 $i$에 의해 생성된다.

\begin{exercise}
\label{fields-exercise-C-not-countably-generated}
$\mathbf{C}$가 $\mathbf{Q}$ 위에서 가산 생성집합을 갖지 않음을 보여라.
\end{exercise}

\noindent
이제 하나의 원소로 생성되는 확대들을 분류하자.

\begin{lemma}[단순 확대의 분류]
\label{fields-lemma-field-extension-generated-by-one-element}
체 확대 $F/k$가 하나의 원소로 생성되면, 이는 유리함수체
$k(t)/k$ 또는 기약인 $P \in k[t]$에 대한 확대
$k[t]/(P)$ 중 하나와 $k$-동형이다.
\end{lemma}

\noindent
가장 중요한 체 확대의 많은 경우가 하나의 원소로 생성됨을 뒤에서
볼 것이므로, 이 결과는 실제로 유용하다.

\begin{proof}
$F = k(\alpha)$인 $\alpha \in F$를 택하자. 가정에 의해 그러한
$\alpha$가 존재한다. 미정원 $t$를 $\alpha$로 보내는 환 사상
$$
k[t] \to F
$$
가 있다. 그 상은 정역이므로 핵은 소 이데알이다. 따라서 핵은
$(0)$이거나 기약인 $P \in k[t]$에 대한 $(P)$이다.

\medskip\noindent
핵이 기약인 $P \in k[t]$에 대한 $(P)$이면, 사상은 $k[t]/(P)$를
통해 분해되고 체의 사상 $k[t]/(P) \to F$를 유도한다. 상이
$\alpha$를 포함하므로 이 사상은 쉽게 전사임을 알 수 있고, 따라서
동형사상이다. 이 경우 $k[t]/(P) \simeq F$이다.

\medskip\noindent
핵이 자명하면 단사사상 $k[t] \to F$를 얻는다. 따라서 분수체
$k(t)$에서 $F$로 가는 사상을 정의할 수 있다. $R(t), Q(t) \in k[t]$인
몫 $R(t)/Q(t)$를 $R(\alpha)/Q(\alpha)$로 보낸다. $k[t] \to F$가
단사라는 가정에 의해 $Q$가 영다항식이 아닌 한
$Q(\alpha) \neq 0$이다. $k[t]$의 분수체는 유리함수체 $k(t)$이므로
상이 $\alpha$를 포함하는 사상 $k(t) \to F$를 얻는다. 따라서 이는
전사이고, 그러므로 동형사상이다.
\end{proof}

\begin{lemma}
\label{fields-lemma-common-extension-field}
$k$를 체라 하고 $E/k$와 $F/k$를 체 확대라 하자. 그러면 공통 체
확대 $M/k$, 즉 $k$의 확대의 사상 $E \to M$과 $F \to M$이 존재하는
확대체가 있다.
\end{lemma}

\begin{proof}
$E$가 $k$의 유한생성 체 확대인 경우만 증명한다. 일반적인 경우는
초른 보조정리 유형의 논증으로 이 경우에서 따른다(세부사항은 생략한다).

\medskip\noindent
먼저 $E$가 $k$의 단순 확대라고 하자. Lemma
\ref{fields-lemma-field-extension-generated-by-one-element}에 의해 이는
$E = k(t)$가 유리함수체인 경우이거나, 어떤 기약다항식
$P \in k[t]$에 대해 $E = k[t]/(P)$인 경우이다. 첫 번째 경우에는
$M = F(t)$를 유리함수체로 택하고 자명한 사상 $E \to M$과
$F \to M$을 취한다. 두 번째 경우에는 $F[t]$에서 $P$의 상의
기약인 인수 $Q$를 택하고, 자명한 사상 $E \to M$과 $F \to M$을
갖는 $M = F[t]/(Q)$를 취한다.

\medskip\noindent
$E = k(\alpha_1, \ldots, \alpha_n)$이면 $n$에 대한 귀납법으로 확대
$M/k$와 사상 $F \to M$ 및
$k(\alpha_1, \ldots, \alpha_{n - 1}) \to M$을 찾을 수 있다.
앞 단락에서 다룬 단순 확대의 경우에 의해 확대
$M'/k(\alpha_1, \ldots, \alpha_{n - 1})$와 사상 $M \to M'$ 및
$k(\alpha_1, \ldots, \alpha_n) \to M'$을 찾을 수 있다.
$M'$을 $k$의 확대라고 보면 원하는 성질을 갖는다.
\end{proof}



\section{유한 확대}
\label{fields-section-finite-extensions}

\noindent
$F/E$가 체 확대이면 $F$는 분명히 $E$ 위의 벡터공간이기도 하다
(스칼라 작용은 단지 $F$에서의 곱셈이다).

\begin{definition}
\label{fields-definition-degree}
$F/E$를 체 확대라 하자. $F$를 $E$-벡터공간으로 보았을 때의 차원을
확대의 {\it 차수}라 하며 $[F : E]$로 나타낸다. $[F : E] < \infty$이면
$F$를 $E$의 {\it 유한} 확대라고 한다.
\end{definition}

\begin{example}
\label{fields-example-C-over-R}
체 $\mathbf{C}$는 기저 $1, i$를 갖는 $\mathbf{R}$ 위의 2차원
벡터공간이다. 따라서 $\mathbf{C}$는 $\mathbf{R}$의 차수 2인 유한
확대이다.
\end{example}

\begin{lemma}
\label{fields-lemma-finite-goes-up}
$K/E/F$를 대수적 체 확대들의 탑이라 하자. $K$가 $F$ 위에서
유한이면 $K$는 $E$ 위에서도 유한이다.
\end{lemma}

\begin{proof}
정의에서 바로 따른다.
\end{proof}

\noindent
이제 가장 중요한 특수한 예, 즉 Lemma
\ref{fields-lemma-field-extension-generated-by-one-element}가 주는 경우의
차수를 다음 두 예에서 살펴보자.

\begin{example}[유리함수체의 차수]
\label{fields-example-degree-rational-function-field}
$k$가 임의의 체이면 유리함수체 $k(t)$는 유한 확대가 {\it 아니다}.
예를 들어 원소들 $\left\{t^n, n \in \mathbf{Z}\right\}$은 $k$ 위에서
선형독립이다.

\medskip\noindent
실제로 $k$가 비가산이면 $k(t)$는 $k$-벡터공간으로서 {\it 비가산}
차원이다. 이를 보이기 위해 원소들의 족
$\{1/(t- \alpha), \alpha \in k\} \subset k(t)$가 $k$ 위에서
선형독립임을 주장한다. 이들 사이의 자명하지 않은 관계는 모순을
낳는다. 예를 들어 $\mathbf{C}$ 위에서 생각하면, 유리형 함수로 본
$\frac{1}{t-\alpha}$는 $\mathbf{C}$의 $\alpha$에서만 극을 가지므로
이를 알 수 있다. 따라서 $c_i \in k^*$이고 $\alpha_i \in k$들이 서로
다를 때 합 $\sum c_i \frac{1}{t - \alpha_i}$는 각 $\alpha_i$에서
극을 갖는다. 특히 영일 수 없다.

\medskip\noindent
흥미롭게도 이는 복소수체 위의 힐베르트 영점정리에 대한 짧은 증명을
준다. 조금 더 일반적인 결과는 Algebra, Theorem
\ref{algebra-theorem-uncountable-nullstellensatz}를 보라.
\end{example}

\begin{lemma}
\label{fields-lemma-finite-finitely-generated}
체의 유한 확대는 유한생성 체 확대이다. 역은 참이 아니다.
\end{lemma}

\begin{proof}
$F/E$를 체의 유한 확대라 하자. $\alpha_1, \ldots, \alpha_n$을
$E$ 위의 벡터공간 $F$의 기저라 하자. 그러면
$F = E(\alpha_1, \ldots, \alpha_n)$이므로 $F/E$는 유한생성 체
확대이다. 역이 참이 아님은
Example \ref{fields-example-degree-rational-function-field}에서 따른다.
\end{proof}

\begin{example}[단순 대수적 확대의 차수]
\label{fields-example-degree-simple-algebraic-extension}
Example \ref{fields-example-monogenic-extension}에서 다룬 형태의 단일생성
체 확대 $E/k$를 생각하자. 다시 말해, $P \in k[t]$가 기약다항식일 때
$E = k[t]/(P)$이다. 그러면 차수 $[E : k]$는 다항식 $P$의 차수
$d = \deg(P)$이다. 실제로
\begin{equation}
\label{fields-equation-P}
P = a_d t^d + a_{d - 1} t^{d - 1} + \ldots + a_0.
\end{equation}
이고 $a_d \not = 0$이라고 하자. $k[t]/(P)$에서
$1, t, \ldots, t^{d - 1}$의 상들은 $k$ 위에서 선형독립이다. 이들
사이의 관계는 $P$보다 차수가 엄밀히 작을 것이지만 $P$가 이데알
$(P)$에서 차수가 가장 작은 원소이기 때문이다.

\medskip\noindent
거꾸로, 집합 $S = \{1, t, \ldots, t^{d - 1}\}$(더 정확히는 그
상들)는 벡터공간 $k[t]/(P)$를 생성한다. 실제로 (\ref{fields-equation-P})에
의해 $a_d t^d$는 $S$의 선형생성에 속한다. $a_d$가 가역이므로
$t^d$도 $S$의 선형생성에 속한다. 마찬가지로 관계 $t P(t) = 0$은
$t^{d + 1}$의 상이 $\{1, t, \ldots, t^d\}$의 선형생성에 속함을
보여 주며, 방금 보인 바에 의해 이는 $S$의 선형생성에 속한다.
위로 귀납하면 $n \geq d$인 $t^n$의 상은 $S$의 선형생성에 속한다.
\end{example}

\noindent
예를 들어 이는 $[\mathbf{C}: \mathbf{R}] = 2$라는 관찰을 확인한다.
더 일반적으로 $k$가 체이고 $\alpha \in k$가 제곱이 아니면 기약다항식
$x^2 - \alpha \in k[x]$를 사용하여 차수 2인 확대
$k[x]/(x^2 - \alpha)$를 구성할 수 있다. 이를
$k(\sqrt{\alpha})$로 쓴다. 명백한 이유로 이러한 확대를
{\it 이차}라고 한다.

\medskip\noindent
차수에 관한 기본 사실은 탑에서 {\it 곱셈적}이라는 것이다.

\begin{lemma}[곱셈성]
\label{fields-lemma-multiplicativity-degrees}
체의 탑 $F/E/k$가 주어졌다고 하자. 그러면
$$
[F:k] = [F:E][E:k]
$$
이다.
\end{lemma}

\begin{proof}
$\alpha_1, \ldots, \alpha_n \in F$를 $F$의 $E$-기저라 하고,
$\beta_1, \ldots, \beta_m \in E$를 $E$의 $k$-기저라 하자. 주장하는
바는 곱들의 집합
$\{\alpha_i \beta_j, 1 \leq i \leq n, 1 \leq j \leq m\}$이
$F$의 $k$-기저라는 것이다. 먼저 이들이 $k$ 위에서 $F$를 생성함을
확인하자.

\medskip\noindent
가정에 의해 $\{\alpha_i\}$는 $E$ 위에서 $F$를 생성한다. 따라서
$f \in F$이면 어떤 $a_i \in E$가 있어
$$
f = \sum\nolimits_i a_i \alpha_i,
$$
이고, 각 $i$에 대하여 어떤 $b_{ij} \in k$를 사용하여
$a_i = \sum b_{ij} \beta_j$로 쓸 수 있다. 이를 합치면
$$
f = \sum\nolimits_{i,j} b_{ij} \alpha_i \beta_j,
$$
이므로 $\{\alpha_i \beta_j\}$가 $k$ 위에서 $F$를 생성한다.

\medskip\noindent
이제 $c_{ij} \in k$에 대해 자명하지 않은 관계
$$
\sum\nolimits_{i,j} c_{ij} \alpha_i \beta_j = 0
$$
가 존재한다고 하자. 그러면
$$
\sum\nolimits_i \alpha_i \left( \sum\nolimits_j c_{ij} \beta_j \right) = 0,
$$
이고 $\beta_j$들이 그러하듯 괄호 안의 항들은 $E$에 속한다.
$\{\alpha_i\}$의 $E$-선형독립성에 의해 괄호 안의 합들은 모두
영이다. 이어서 $\{\beta_j\}$의 $k$-선형독립성에 의해 모든
$c_{ij}$가 영이다.
\end{proof}

\noindent
잠시 곁가지인 정의를 하나 들자.

\begin{definition}
\label{fields-definition-number-field}
체 $K$의 표수가 $0$이고 확대 $K/\mathbf{Q}$가 유한이면 그 체를
{\it 수체}라고 한다.
\end{definition}

\noindent
수체는 대수적 수론의 기본 대상이다. 뒤에서 수체 안의 정수
$\mathbf{Z}$의 유사물에 대해서도 일종의 유일인수분해가 여전히
성립함을 볼 것이다(엄밀한 유일인수분해는 일반적으로 성립하지
않는다!).


\section{대수적 확대}
\label{fields-section-algebraic-extensions}

\noindent
중요한 확대의 한 부류는 모든 원소가 유한 확대를 생성하는 경우이다.
% P01-E0239: frozen subject--verb disagreement recorded in the sidecar.

\begin{definition}
\label{fields-definition-algebraic}
체 확대 $F/E$를 생각하자. 원소 $\alpha \in F$가 계수가 $E$에 있는
어떤 영이 아닌 다항식의 근이면 $\alpha$가 $E$ 위에서
{\it 대수적}이라고 한다. $F$의 모든 원소가 대수적이면 $F$를
$E$의 {\it 대수적 확대}라고 한다.
\end{definition}

\noindent
Lemma \ref{fields-lemma-field-extension-generated-by-one-element}에 의해
부분확대 $E(\alpha)$는 유리함수체 $E(t)$ 또는 기약다항식
$P \in E[t]$에 대한 몫환 $E[t]/(P)$와 동형이다. 뒤의 경우에는
$\alpha$가 $E$ 위에서 대수적이고(실제로 Lemma
\ref{fields-lemma-field-extension-generated-by-one-element}의 증명은
$\alpha$가 $P$의 근이 되도록 $P$를 택할 수 있음을 보여 준다),
앞의 경우에는 대수적이지 않다.

\begin{example}
\label{fields-example-C-algebraic-over-R}
체 $\mathbf{C}$는 $\mathbf{R}$ 위에서 대수적이다. 실제로
$\mathbf{C}$의 $\alpha = a + ib$에 대하여
$\alpha^2 - 2a\alpha + a^2 + b^2 = 0$은 $\alpha$가
$\mathbf{R}$ 위에서 만족하는 다항식 방정식이다.
\end{example}

\begin{example}
\label{fields-example-compact-riemann-surface-is-finite-over-P1}
$X$를 콤팩트 리만 곡면이라 하고,
$f \in \mathbf{C}(X) - \mathbf{C}$를 $X$ 위의 임의의 비상수 유리형
함수라 하자(Example \ref{fields-example-field-of-meromorphic-functions}를
보라). 그러면 $\mathbf{C}(X)$가 $f$에 의해 생성되는 부분확대
$\mathbf{C}(f)$ 위에서 대수적임이 알려져 있다. 여기서는 증명하지 않는다.
% P01-E1335: frozen example setup omits “be” after the displayed element.
\end{example}

\begin{lemma}
\label{fields-lemma-algebraic-goes-up}
$K/E/F$를 체 확대들의 탑이라 하자.
\begin{enumerate}
\item $\alpha \in K$가 $F$ 위에서 대수적이면 $\alpha$는 $E$ 위에서도
대수적이다.
\item $K$가 $F$ 위에서 대수적이면 $K$는 $E$ 위에서도 대수적이다.
\end{enumerate}
\end{lemma}

\begin{proof}
정의에서 바로 따른다.
\end{proof}

\noindent
이제 유한성과 대수성 사이에 깊은 관계가 있음을 보이자.

\begin{lemma}
\label{fields-lemma-finite-is-algebraic}
유한 확대는 대수적이다. 실제로 확대 $E/k$가 대수적일 필요충분조건은
어떤 $\alpha \in E$가 생성하는 모든 부분확대 $k(\alpha)/k$가
유한인 것이다.
\end{lemma}

\noindent
일반적으로 대수적 확대가 유한이라는 명제는 전혀 참이 아니다.

\begin{proof}
$E/k$를 차수가 $n$인 유한 확대라 하고 $\alpha \in E$를 택하자.
원소 $1, \alpha, \ldots, \alpha^n$은 $k$ 위에서 선형종속이다.
그렇지 않으면 반드시 $[E : k] > n$이기 때문이다. 이제 하나의
선형종속 관계가 $\alpha$가 만족해야 하는 다항식을 준다.

\medskip\noindent
마지막 명제에 대해서는, Examples
\ref{fields-example-degree-rational-function-field}와
\ref{fields-example-degree-simple-algebraic-extension}에 의해 단일생성 확대
$k(\alpha)/k$가 유한일 필요충분조건이 $\alpha$가 $k$ 위에서
대수적인 것임에 유의하자. 따라서 $E/k$가 대수적이면 각
$k(\alpha)/k$, $\alpha \in E$는 유한 확대이고, 그 역도 성립한다.
\end{proof}

\noindent
앞의 증명(실제로는 Examples
\ref{fields-example-degree-rational-function-field}와
\ref{fields-example-degree-simple-algebraic-extension})에서 다음 보조정리를
추출할 수 있다. 단일생성 확대가 유한일 필요충분조건은 그것이
대수적인 것이다. 다음 결과에서 이 관찰을 사용한다.

\begin{lemma}
\label{fields-lemma-algebraic-finitely-generated}
\begin{slogan}
유한생성 대수적 확대는 유한이다.
\end{slogan}
$k$를 체라 하고 $\alpha_1, \alpha_2, \ldots, \alpha_n$을 어떤
확대체의 원소들이라 하자. 각 $\alpha_i$가 $k$ 위에서 대수적이면
확대 $k(\alpha_1, \ldots, \alpha_n)/k$는 유한이다. 즉 유한생성
대수적 확대는 유한이다.
\end{lemma}

\begin{proof}
실제로 각 확대
$k(\alpha_{1}, \ldots, \alpha_{i+1})/k(\alpha_1, \ldots, \alpha_{i})$
는 한 원소로 생성되고 대수적이므로 유한이다. 차수의 곱셈성
(Lemma \ref{fields-lemma-multiplicativity-degrees})에 의해 결론을 얻는다.
\end{proof}

\noindent
$\mathbf{Q}$ 위에서 대수적인 복소수들을 단순히 {\it 대수적 수}라고
한다. 예를 들어 $\sqrt{2}$는 대수적이고 $i$도 대수적이지만 $\pi$는
그렇지 않다. 대수적 수들이 체를 이룬다는 것은 기본적인 사실이다.
그러나 수가 유리수 계수의 영이 아닌 다항식 방정식을 만족할 때
정확히 대수적이라는 정의만으로부터 이를 어떻게 증명할지는
명백하지 않다(예를 들어 다항식 방정식을 이용해서 말이다).
% P01-E0240: frozen subject--verb disagreement recorded in the sidecar.

\begin{lemma}
\label{fields-lemma-algebraic-elements}
$E/k$를 체 확대라 하자. $E$에서 $k$ 위에 대수적인 원소들은
$E/k$의 부분확대를 이룬다.
\end{lemma}

\begin{proof}
$\alpha, \beta \in E$가 $k$ 위에서 대수적이라 하자. Lemma
\ref{fields-lemma-algebraic-finitely-generated}에 의해
$k(\alpha, \beta)/k$는 유한 확대이다. 따라서
$k(\alpha + \beta) \subset k(\alpha, \beta)$는 유한 확대이고,
Lemma \ref{fields-lemma-finite-is-algebraic}에 의해 $\alpha + \beta$는
대수적이다. $\alpha$와 $\beta$의 차, 곱, 몫도 마찬가지이다.
\end{proof}

\noindent
환의 경우와 마찬가지로, 체 확대의 여러 좋은 성질은 체의 탑과
합성체에 의해 보존된다.

\begin{lemma}
\label{fields-lemma-algebraic-permanence}
$E/k$와 $F/E$를 체의 대수적 확대라 하자. 그러면 $F/k$는 체의
대수적 확대이다.
\end{lemma}

\begin{proof}
$\alpha \in F$를 택하자. 그러면 $\alpha$는 $E$ 위에서 대수적이다.
핵심 관찰은 $\alpha$가 $k$의 어떤 유한생성 확대 위에서 대수적이라는
것이다.
% P01-E0241: frozen containment direction recorded in the sidecar.
즉, $\alpha $가 $k(S)$ 위에서 대수적이 되게 하는 유한집합
$S \subset E$가 존재한다. 실제로 대수적이라는 것은 $\alpha$가
만족하는 어떤 다항식이 $E[x]$에 존재한다는 뜻이고, 그 다항식의
계수들을 $S$로 택할 수 있으므로 이는 명백하다. 따라서 $\alpha$는
$k(S)$ 위에서 대수적이다. 특히 확대 $k(S, \alpha)/ k(S)$는
유한이다. $S$가 유한집합이고 $k(S)/k$가 대수적이므로 Lemma
\ref{fields-lemma-algebraic-finitely-generated}에 의해 $k(S)/k$는 유한이다.
곱셈성(Lemma \ref{fields-lemma-multiplicativity-degrees})을 사용하면
$k(S,\alpha)/k$가 유한임을 얻으므로 $\alpha$는 $k$ 위에서 대수적이다.
\end{proof}

\noindent
앞의 논증에서 사용한 증명 방법, 즉 $E$ 위에서 대수적이라는 성질이
$E$의 유한생성 부분확대로 {\it 하강했다}는 생각은 대수학 전체에서
거듭 나타난다. 예를 들어 이를 통해 일반적인 가환대수 문제를
뇌터인 경우로 환원할 수 있는 일이 많다.

\begin{lemma}
\label{fields-lemma-size-algebraic-extension}
$E/F$를 체의 대수적 확대라 하자. 그러면 $E$의 기수 $|E|$는
기껏해야 $\max(\aleph_0, |F|)$이다.
\end{lemma}

\begin{proof}
$S$를 계수가 $F$에 있는 비상수 다항식들의 집합이라 하자. 각
$P \in S$에 대하여 근들의 집합
$r(P, E) = \{\alpha \in E \mid P(\alpha) = 0\}$은 유한이다
(세부사항은 생략한다). 또한 $E$가 $F$ 위에서 대수적이므로
$E = \bigcup_{P \in S} r(P, E)$이다. $S$는 $F$의 복사본들의
유한곱들의 가산 합집합이므로 그 기수가 $\max(\aleph_0, |F|)$로
제한됨은 명백하다. 따라서 $E$의 기수도 그러하다.
\end{proof}

\begin{lemma}
\label{fields-lemma-subalgebra-algebraic-extension-field}
$E/F$를 유한 확대, 더 일반적으로는 체의 대수적 확대라 하자.
임의의 부분환 $F \subset R \subset E$는 체이다.
\end{lemma}

\begin{proof}
$\alpha \in R$을 영이 아닌 원소라 하자. 그러면
$1, \alpha, \alpha^2, \ldots$는 $R$에 속한다. Lemma
\ref{fields-lemma-finite-is-algebraic}에 의해 $a_i \in F$인 자명하지 않은 관계
$a_0 + a_1 \alpha + \ldots + a_d \alpha^d = 0$를 얻는다.
$a_0 \not = 0$이라고 가정해도 된다. 그렇지 않으면 관계를
$\alpha$로 나누어 $d$를 줄일 수 있기 때문이다. 그러면
$$
a_0 = \alpha (- a_1  - \ldots - a_d \alpha^{d - 1})
$$
이고, 이는 $\alpha$의 역원이 $R$의 원소
$a_0^{-1} (- a_1  - \ldots - a_d \alpha^{d - 1})$임을 증명한다.
\end{proof}

\begin{lemma}
\label{fields-lemma-algebraic-extension-self-map}
$E/F$를 체의 대수적 확대라 하자. 임의의 $F$-대수 사상
$f : E \to E$는 자기동형사상이다.
% P01-E0242: frozen lemma setup omits “be”.
\end{lemma}

\begin{proof}
$E/F$가 유한이면 $f : E \to E$는 유한차원 벡터공간들의
$F$-선형 단사사상(Lemma \ref{fields-lemma-field-maps-injective})이므로
전단사이다. 일반적인 경우에도 $f$가 단사임을 알 수 있다.
$\alpha \in E$라 하고 $P \in F[x]$를 $P(\alpha) = 0$인 다항식이라
하자. $E' \subset E$를 $P$의 $E$ 안의 근들
$\alpha = \alpha_1, \ldots, \alpha_n$이 생성하는 $E$의 부분체라
하자. Lemma \ref{fields-lemma-algebraic-finitely-generated}에 의해 $E'$은
$F$ 위에서 유한이다. $f$가 근들의 집합을 보존하므로
$f|_{E'} : E' \to E'$임을 얻는다. 따라서 증명의 첫 부분에 의해
$f|_{E'}$는 동형사상이고, $\alpha$가 $f$의 상에 속한다는 결론을
얻는다.
\end{proof}





\section{최소다항식}
\label{fields-section-minimal-polynomials}

\noindent
$E/k$를 체 확대라 하고 $\alpha \in E$가 $k$ 위에서 대수적이라
하자. 그러면 $\alpha$는 $k[x]$ 안의 (자명하지 않은) 다항식 방정식을
만족한다. $P(\alpha) = 0$인 다항식 $P \in k[x]$들의 집합을 생각하자.
가정에 의해 이 집합은 영다항식만 포함하지는 않는다. 이 집합이
{\it 이데알}임은 쉽게 알 수 있다. 실제로 이는 사상
$$
k[x] \to E, \quad x \mapsto \alpha
$$
의 핵이다. $k[x]$는 주 아이디얼 정역이므로 이 이데알의
{\it 생성원} $P \in k[x]$가 존재한다. 일반성을 잃지 않고 $P$가
모닉이라고 가정하면 $P$는 유일하게 결정된다.

\begin{definition}
\label{fields-definition-minimal-polynomial}
위의 다항식 $P$를 $\alpha$의 $k$ 위 {\it 최소다항식}이라 한다.
\end{definition}

\noindent
최소다항식은 다음과 같이 특징지어진다. 이는 $\alpha$를 영이 되게 하는
가장 낮은 차수의 모닉 다항식이다. $P$의 임의의 비상수 배수는 더 높은
차수를 가지며, $P$의 배수만이 $\alpha$를 영이 되게 할 수 있다.
이것이 {\it 최소}라는 이름을 설명한다.

\medskip\noindent
최소다항식은 명백히 {\it 기약}이다. 이는 $k[x]$에서 $\alpha$를
영이 되게 하는 다항식들로 이루어진 이데알이 소 이데알이라는 명제와
동치이다. 사상 $k[x] \to E, x \mapsto \alpha$는 정역(심지어 체)으로
가는 사상이므로 그 핵이 소 이데알이라는 사실에서 이를 얻는다.

\begin{lemma}
\label{fields-lemma-degree-minimal-polynomial}
최소다항식의 차수는 $[k(\alpha) : k]$이다.
\end{lemma}

\begin{proof}
이는 Lemma \ref{fields-lemma-field-extension-generated-by-one-element}의 논증을
다시 말한 것일 뿐이다. $P$가 $\alpha$의 최소다항식이면 앞서 언급한
증명에서와 같이 사상
$$
k[x]/(P) \to k(\alpha), \quad x \mapsto \alpha
$$
는 동형사상이고, 그러한 확대의 차수는 이미 계산했다(Example
\ref{fields-example-degree-simple-algebraic-extension}를 보라).
\end{proof}

\noindent
따라서 위 증명의 관찰은 $\alpha \in E$가 대수적이면
$k(\alpha) \subset E$가 $k[x]/(P)$와 동형이라는 것이다.


\section{대수적 폐포}
\label{fields-section-algebraic-closure}

\noindent
``대수학의 기본정리''는 $\mathbf{C}$가 대수적으로 닫혀 있다고
말한다. 이 결과의 아름다운 증명은 복소해석학의 리우빌 정리를 사용한다.
여기서는 다른 증명을 제시한다(Lemma
\ref{fields-lemma-C-algebraically-closed}를 보라).
% P01-E0243: frozen comma splice recorded in the sidecar.

\begin{definition}
\label{fields-definition-algebraically-closed}
체 $F$의 모든 대수적 확대 $E/F$가 자명한, 즉 $E = F$인 경우
그 체를 {\it 대수적으로 닫힌 체}라고 한다.
\end{definition}

\noindent
모든 문헌에서 이 정의를 쓰는 것은 아닐 수 있다. 다음 보조정리는
이 정의를 다른 정의와 비교한다.

\begin{lemma}
\label{fields-lemma-algebraically-closed}
$F$를 체라 하자. 다음 조건들은 동치이다.
\begin{enumerate}
\item $F$는 대수적으로 닫혀 있다.
\item $F$ 위의 모든 기약다항식은 일차이다.
\item $F$ 위의 모든 비상수 다항식은 근을 갖는다.
\item $F$ 위의 모든 비상수 다항식은 일차 인수들의 곱이다.
\end{enumerate}
\end{lemma}

\begin{proof}
$F$가 대수적으로 닫혀 있으면 모든 기약다항식은 일차이다. 실제로
차수가 $> 1$인 기약다항식이 존재하면 이는 자명하지 않은 유한
(따라서 대수적) 체 확대를 생성한다. Example
\ref{fields-example-degree-simple-algebraic-extension}를 보라. 따라서 (1)은
(2)를 함의한다. 모든 기약다항식이 일차이면 모든 기약다항식은 근을
가지므로 모든 비상수 다항식이 근을 갖는다. 따라서 (2)는 (3)을
함의한다.

\medskip\noindent
모든 비상수 다항식이 근을 갖는다고 하자. $P \in F[x]$를
비상수라 하자. $P(\alpha) = 0$이고 $\alpha \in F$이면 나머지 있는
나눗셈에 의해 어떤 $Q \in F[x]$에 대하여 $P = (x - \alpha)Q$이다.
따라서 차수에 대한 귀납법으로 임의의 비상수 다항식을
$c \prod (x - \alpha_i)$라는 곱으로 쓸 수 있다.

\medskip\noindent
마지막으로 $F$ 위의 모든 비상수 다항식이 일차 인수들의 곱이라고
하자. $E/F$를 대수적 확대라 하자. 그러면 $E$의 모든 단순 부분확대
$F(\alpha)/F$는 반드시 자명하다(가정에 의해 기약다항식은 일차뿐이기
때문이다). 따라서 $E = F$이다. (4)가 (1)을 함의함을 보였으므로
증명이 끝났다.
\end{proof}

\noindent
이제 체의 ``보편적인'' 대수적 확대를 정의하고자 한다. 실제로는
주의해야 한다. 대수적 폐포는 보편 대상이 {\it 아니다}. 즉 대수적
폐포는 {\it 유일한} 동형사상을 제외하고 유일한 것이 아니라, 단지
동형사상을 제외하고 유일할 뿐이다. 그럼에도 함자적이지는 않지만
매우 편리하다.

\begin{definition}
\label{fields-definition-algebraic-closure}
$F$를 체라 하자. $F$의 {\it 대수적 폐포}란 $F$를 포함하고 다음을
만족하는 체 $\overline{F}$이다.
\begin{enumerate}
\item $\overline{F}$는 $F$ 위에서 대수적이다.
\item $\overline{F}$는 대수적으로 닫혀 있다.
\end{enumerate}
\end{definition}

\noindent
$F$가 대수적으로 닫혀 있으면 $F$는 자신의 대수적 폐포이다.
이제 기본적인 존재 결과를 증명한다.

\begin{theorem}
\label{fields-theorem-existence-algebraic-closure}
모든 체는 대수적 폐포를 갖는다.
\end{theorem}

\noindent
이 증명은 이 장의 나머지 부분에는 대체로 곁가지이다. 그러나 체를
대수적으로 닫힌 체에 매입하는 것이 {\it 가능하다}는 사실은 알아야
하며, 우리는 흔히 그러한 매입이 이미 이루어졌다고 가정할 것이다.

\begin{proof}
$F$를 체라 하자. Lemma \ref{fields-lemma-size-algebraic-extension}에 의해
$F$의 대수적 확대의 기수는 $\max(\aleph_0, |F|)$로 제한된다.
$F$를 포함하면서 $|S| > \max(\aleph_0, |F|)$인 집합 $S$를 택하자.
다음 삼중항 $(E, \sigma_E, \mu_E)$들을 생각한다.
\begin{enumerate}
\item $E$는 $F \subset E \subset S$인 집합이다.
\item $\sigma_E : E \times E \to E$와 $\mu_E : E \times E \to E$는
$(E, \sigma_E, \mu_E)$가 $F$의 체 확대 구조를 정의하게 하는 집합의
사상들이다(특히 $a, b \in F$에 대하여
$\sigma_E(a, b) = a +_F b$이고 $\mu_E$에 대해서도 마찬가지이다).
\item $E/F$는 대수적 체 확대이다.
\end{enumerate}
모든 삼중항 $(E, \sigma_E, \mu_E)$들의 모임은 집합 $I$를 이룬다.
$i \in I$에 대하여 대응하는 $F$의 체 확대를
$E_i = (E_i, \sigma_i, \mu_i)$로 나타낸다.
% P01-E0244: frozen denote-construction and extension direction recorded.
$I$에 다음과 같이 부분순서를 정의한다. $i \leq i'$일 필요충분조건은
$E_i \subset E_{i'}$이고(이는 $E_i$와 $E_{i'}$가 같은 집합 $S$의
부분집합이므로 의미가 있다),
$\sigma_i = \sigma_{i'}|_{E_i \times E_i}$ 및
$\mu_i = \mu_{i'}|_{E_i \times E_i}$인 것이다. 다시 말해 $E_{i'}$는
$E_i$의 체 확대이다.

\medskip\noindent
$T \subset I$를 전순서 부분집합이라 하자. 그러면 유도된 사상
$\sigma_T = \bigcup \sigma_i$와 $\mu_T = \bigcup \mu_i$를 갖춘
$E_T = \bigcup_{i \in T} E_i$가 $I$의 또 다른 원소임은 명백하다.
다시 말해 $I$의 모든 전순서 부분집합은 $I$ 안에서 상계를 갖는다.
% P01-E0245: frozen “totally order” and “a upper bound” forms recorded.
초른 보조정리에 의해 $I$에는 극대원 $(E, \sigma_E, \mu_E)$가
존재한다. $E$가 대수적 폐포라고 주장한다. $I$의 정의에 의해 확대
$E/F$가 대수적이므로 $E$가 대수적으로 닫혀 있음을 보이면 충분하다.

\medskip\noindent
이를 보이기 위해 모순을 이끌어 낸다. $E$가 대수적으로 닫혀 있지
않다고 하자. 그러면 Lemma \ref{fields-lemma-algebraically-closed}에 의해
$E$ 위에 차수가 $> 1$인 기약다항식 $P$가 존재한다. Lemma
\ref{fields-lemma-finite-is-algebraic}에 의해 자명하지 않은 유한 확대
$E' = E[x]/(P)$를 얻는다. Lemma
\ref{fields-lemma-algebraic-permanence}에 의해 $E'/F$가 대수적임에 유의하자.
따라서 $E'$의 기수는 $\leq \max(\aleph_0, |F|)$이다. 초등적인
집합론에 의해 주어진 단사 $E \subset S$를 단사 $E' \to S$로
확장할 수 있다. 다시 말해 $E'$을 우리 집합 $I$의 원소로 볼 수
있는데, 이는 $E$의 극대성과 모순이다. 이 모순으로 증명이 끝난다.
\end{proof}

\begin{lemma}
\label{fields-lemma-map-into-algebraic-closure}
$F$를 체라 하자. $\overline{F}$를 $F$의 대수적 폐포라 하고,
$M/F$를 대수적 확대라 하자. 그러면 $F$-확대의 사상
$M \to \overline{F}$가 존재한다.
\end{lemma}

\begin{proof}
$F \subset E \subset M$가 부분확대이고
$\varphi : E \to \overline{F}$가 $F$-확대의 사상인 순서쌍
$(E, \varphi)$들의 집합을 $I$라 하자. $E \subset E'$이고
$\varphi'|_E = \varphi$일 필요충분조건으로
$(E, \varphi) \leq (E', \varphi')$를 정의하여 집합 $I$에 부분순서를
준다. $T = \{(E_t,  \varphi_t)\} \subset I$가 전순서 부분집합이면
$\bigcup \varphi_t : \bigcup E_t \to \overline{F}$는 $I$의 원소이다.
따라서 $I$의 모든 전순서 부분집합은 상계를 갖는다. 초른 보조정리에
의해 $I$에는 극대원 $(E, \varphi)$가 존재한다. $E = M$이라고
주장한다. 그러면 증명이 끝난다. 그렇지 않으면
$\alpha \in M$, $\alpha \not \in E$를 택하자. 그러면 $\alpha$는
$E$ 위에서 대수적이다. Lemma \ref{fields-lemma-algebraic-goes-up}를 보라.
% P01-E0246: frozen “The alpha” typo recorded in the sidecar.
$P$를 $\alpha$의 $E$ 위 최소다항식이라 하자. $P^\varphi$를
$\varphi$에 의한 $P$의 $\overline{F}[x]$ 안의 상이라 하자.
$\overline{F}$가 대수적으로 닫혀 있으므로 $P^\varphi$는
$\overline{F}$ 안에 근 $\beta$를 갖는다. 그러면 $x$를 $\beta$로
보내어 $\varphi$를
$\varphi' : E(\alpha) = E[x]/(P) \to \overline{F}$로 확장할 수 있다.
이는 원하는 대로 $(E, \varphi)$의 극대성과 모순이다.
\end{proof}

\begin{lemma}
\label{fields-lemma-algebraic-closures-isomorphic}
한 체의 임의의 두 대수적 폐포는 동형이다.
\end{lemma}

\begin{proof}
$F$를 체라 하자. $M$과 $\overline{F}$가 $F$의 대수적 폐포이면
Lemma \ref{fields-lemma-map-into-algebraic-closure}에 의해 $F$-확대의 사상
$\varphi : M \to \overline{F}$가 존재한다. 이제 상 $\varphi(M)$은
대수적으로 닫혀 있다. 한편 Lemma \ref{fields-lemma-algebraic-goes-up}에 의해
확대 $\varphi(M) \subset \overline{F}$는 대수적이다. 따라서
$\varphi(M) = \overline{F}$이다.
\end{proof}





\section{서로소 다항식}
\label{fields-section-relatively-prime}

\noindent
$K$를 대수적으로 닫힌 체라 하자. 그러면 Lemma
\ref{fields-lemma-algebraically-closed}에서 본 것처럼 환 $K[x]$의 이데알
구조는 매우 단순하다. 특히 모든 다항식 $P \in K[x]$는
$$
P = c(x - \alpha_1) \ldots (x - \alpha_n),
$$
로 쓸 수 있다. 여기서 $c$는 상수항이고
$\alpha_1, \ldots, \alpha_n \in k$는 $P$의 근들이다(중복도까지 세었다).
% P01-E0247: frozen text calls c the constant rather than leading coefficient.
% P01-E0248: frozen text places the roots in undefined lowercase k.
명백히 $K[x]$의 기약다항식은 일차다항식 $c(x - \alpha)$,
$c, \alpha \in K$뿐이다(그리고 $c \neq 0$이다).

\begin{definition}
\label{fields-definition-relatively-prime}
$k$가 임의의 체일 때 $k[x]$ 안의 두 다항식이 $k[x]$의 단위 이데알을
생성하면 두 다항식이 {\it 서로소}라고 한다.
\end{definition}

\noindent
위의 논의를 계속하자. $K$가 대수적으로 닫힌 체이면 $K[x]$ 안의
두 다항식이 서로소일 필요충분조건은 공통근이 없는 것이다. 이는
$K[x]$의 극대 이데알들이 $(x - \alpha)$, $\alpha \in K$의 꼴이기
때문이다. 따라서 $F, G \in K[x]$가 공통근을 갖지 않으면 $(F, G)$는
어떤 $(x - \alpha)$에도 포함될 수 없다(포함되면 $\alpha$가
공통근이 될 것이기 때문이다).

\medskip\noindent
$k$가 대수적으로 닫힌 체가 {\it 아니어도}, 위의 논의는 $k[x]$ 안의
두 다항식이 언제 단위 이데알을 생성하는지에 관한 정보를 준다.

\begin{lemma}
\label{fields-lemma-relatively-prime-polynomials}
$k[x]$ 안의 두 다항식이 서로소일 필요충분조건은 $k$의 대수적 폐포
$\overline{k}$에서 공통근을 갖지 않는 것이다.
\end{lemma}

\begin{proof}
주장은 임의의 두 다항식 $P, Q$가 $k[x]$에서 $(1)$을 생성할
필요충분조건이 $\overline{k}[x]$에서 $(1)$을 생성하는 것이라는
것이다. 이는 선형대수학의 한 사실이다. 계수가 $k$에 있는 연립
일차방정식이 해를 가질 필요충분조건은 $k$의 임의의 확대에서 해를
갖는 것이다. 따라서 대수적으로 닫힌 체의 경우로 환원할 수 있고,
그 경우에는 이미 증명한 사실에서 결과가 명백하다.
\end{proof}





\section{분리 대수적 확대}
\label{fields-section-separable-extensions}

\noindent
표수 $p$에서는 체 위의 기약다항식에 특이한 일이 일어난다. 다음
보조정리에서 이를 설명한다.

\begin{lemma}
\label{fields-lemma-irreducible-polynomials}
$F$를 체라 하고 $P \in F[x]$를 $F$ 위의 기약다항식이라 하자.
$P' = \text{d}P/\text{d}x$를 $P$의 $x$에 대한 도함수라 하자.
그러면 다음 두 경우 중 하나가 성립한다.
\begin{enumerate}
\item $P$와 $P'$가 서로소이다. 또는
\item $P'$가 영다항식이다.
\end{enumerate}
두 번째 경우는 $F$의 표수가 $p > 0$일 때만 일어날 수 있다.
이 경우 $q = p^f$가 $p$의 거듭제곱이고 $Q \in F[x]$가 $Q$와
$Q'$가 서로소인 기약다항식일 때 $P(x) = Q(x^q)$이다.
\end{lemma}

\begin{proof}
$P'$의 차수는 $< \deg(P)$임에 유의하자. 따라서 $P$와 $P'$가
서로소가 아니면 $(P, P') = (R)$이고, 여기서 $R$은 차수가
$< \deg(P)$인 다항식이어서 $P$의 기약성과 모순이다. 이로써
(1)과 (2)의 양자택일을 얻는다.

\medskip\noindent
(2)의 경우이고 $P = a_d x^d + \ldots + a_0$이라고 하자. 그러면
$P' = da_d x^{d - 1} + \ldots + a_1$이다. 표수 $0$에서는 이것이
$a_d, \ldots, a_1 = 0$을 강제하여 $P$가 상수가 되므로 모순이다.
따라서 표수 $p$가 양수라는 결론을 얻는다. 이 경우 조건 $P' = 0$은
$p$가 $i$를 나누지 않을 때마다 $a_i = 0$을 강제한다. 다시 말해,
어떤 비상수 다항식 $P_1$에 대하여 $P(x) = P_1(x^p)$이다. 명백히
$P_1$도 기약이다. 차수에 대한 귀납법에 의해 보조정리의 명제와 같이
$P_1(x) = Q(x^q)$임을 알 수 있고, 따라서 $P(x) = Q(x^{pq})$이다.
보조정리가 증명되었다.
\end{proof}

\begin{definition}
\label{fields-definition-separable}
$F$를 체라 하고 $K/F$를 체 확대라 하자.
\begin{enumerate}
\item $F$ 위의 기약다항식 $P$가 그 도함수와 서로소이면 그것이
{\it 분리 가능}하다고 한다.
\item $F$ 위에서 대수적인 $\alpha \in K$가 주어졌을 때, 그
최소다항식이 $F$ 위에서 분리 가능하면 $\alpha$가 $F$ 위에서
{\it 분리 가능}하다고 한다.
\item $K$가 $F$의 대수적 확대일 때 $K$의 모든 원소가 $F$ 위에서
분리 가능하면 $K$가 $F$ 위에서 {\it 분리 가능}\footnote{대수적이지
않은 확대에 대해서는 이 정의가 의미를 갖지 않으며 올바른 정의도
아니다. 독자는 Algebra, Sections \ref{algebra-section-separability}와
\ref{algebra-section-separability-continued}를 보라.}하다고 한다.
\end{enumerate}
\end{definition}

\noindent
Lemma \ref{fields-lemma-irreducible-polynomials}에 의해 표수 $0$에서는 모든
기약다항식이 분리 가능하고, 확대의 모든 대수적 원소가 분리 가능하며,
모든 대수적 확대가 분리 가능하다.

\begin{lemma}
\label{fields-lemma-separable-goes-up}
$K/E/F$를 대수적 체 확대들의 탑이라 하자.
\begin{enumerate}
\item $\alpha \in K$가 $F$ 위에서 분리 가능하면 $\alpha$는 $E$
위에서도 분리 가능하다.
\item $K$가 $F$ 위에서 분리 가능하면 $K$는 $E$ 위에서도 분리
가능하다.
\end{enumerate}
\end{lemma}

\begin{proof}
이하에서 Lemma \ref{fields-lemma-irreducible-polynomials}를 따로 언급하지
않고 사용한다. $P$를 $\alpha$의 $F$ 위 최소다항식이라 하고,
$Q$를 $\alpha$의 $E$ 위 최소다항식이라 하자. 그러면 $Q$는 다항식환
$E[x]$에서 $P$를 나누므로 $P = QR$이라고 쓸 수 있다. 이때
$P' = Q'R + QR'$이다. 따라서 $Q' = 0$이면 $Q$가 $P$와 $P'$를
모두 나누고, 보조정리에 의해 $P' = 0$이다. 이는 (1)을 증명한다.
(2)는 (1)과 정의에서 바로 따른다.
\end{proof}

\begin{lemma}
\label{fields-lemma-recognize-separable}
$F$를 체라 하자. $F$ 위의 기약다항식 $P$가 분리 가능할
필요충분조건은 $F$의 대수적 폐포에서 $P$가 서로 다른 근들을 갖는
것이다.
\end{lemma}

\begin{proof}
$\alpha \in \overline{F}$가 $P$와 $P'$의 공통근이라고 하자.
그러면 어떤 다항식 $Q$에 대하여 $P = (x - \alpha)Q$이다. 미분하면
$P' = Q + (x - \alpha)Q'$을 얻는다. 따라서 $\alpha$는 $Q$의
근이다. 그러므로 $P$와 $P'$가 공통근을 가지면 $P$는 서로 다른
근들만 갖지 않는다. 거꾸로 $P$가 중근을 갖는, 즉
$(x - \alpha)^2$이 $P$를 나누는 경우 $\alpha$는 $P$와 $P'$의
공통근이다. Lemma \ref{fields-lemma-relatively-prime-polynomials}와 결합하면
보조정리가 증명된다.
\end{proof}

\begin{lemma}
\label{fields-lemma-nr-roots-unchanged}
$F$를 체라 하고 $\overline{F}$를 $F$의 대수적 폐포라 하자. $F$의
표수를 $p > 0$이라 하고 $P$를 $F$ 위의 다항식이라 하자. 그러면
$\overline{F}$에서 $P$의 근들의 집합과 $P(x^p)$의 근들의 집합은
같은 기수를 갖는다(중복도는 세지 않는다).
% P01-E0249: frozen subject--verb disagreement recorded in the sidecar.
\end{lemma}

\begin{proof}
명백히 $\alpha$가 $P(x^p)$의 근일 필요충분조건은 $\alpha^p$가
$P$의 근인 것이다. 다시 말해 $P(x^p)$의 근들은, $\beta$가 $P$의
근일 때 $x^p - \beta$의 근들이다. 따라서 사상
$\overline{F} \to \overline{F}$, $\alpha \mapsto \alpha^p$가
전단사임을 보이면 충분하다. $\overline{F}$가 대수적으로 닫혀 있어서
모든 원소가 $p$제곱근을 가지므로 이 사상은 전사이다. 표수가 $p$이므로
$\alpha^p = \beta^p$이면 $(\alpha - \beta)^p = 0$이고, 따라서
단사이다. 물론 체에서는 $x^p = 0$이면 $x = 0$이다.
\end{proof}

\noindent
$F$를 체라 하고 $P$를 $F$ 위의 기약다항식이라 하자. 그러면 어떤
분리 가능한 기약다항식 $Q$에 대하여 $P = Q(x^q)$임을 안다(Lemma
\ref{fields-lemma-irreducible-polynomials}). 여기서 $q$는 표수 $p$의
거듭제곱이다(표수가 영이면 $q = 1$\footnote{이 장에서는
$0^0 = 1$로 두는 것이 좋은 관례이다.}이고 $Q = P$이다). Lemma
\ref{fields-lemma-nr-roots-unchanged}에 의해 $F$의 임의의 대수적 폐포에서
$P$와 $Q$의 근의 수는 같다. Lemma \ref{fields-lemma-recognize-separable}에
의해 이 수는 $Q$의 차수와 같다.

\begin{definition}
\label{fields-definition-separable-degree}
$F$를 체라 하고 $P$를 $F$ 위의 기약다항식이라 하자. $P$의
{\it 분리차수}란 $F$의 임의의 대수적 폐포에서 $P$의 근들의 집합이
갖는 기수이다(위의 논의를 보라). 기호는 $\deg_s(P)$이다.
\end{definition}

\noindent
$P$의 분리차수는 항상 그 차수를 나누며, 그 몫은 표수의 거듭제곱이다.
표수가 영이면 $\deg_s(P) = \deg(P)$이다.

\begin{situation}
\label{fields-situation-finitely-generated}
여기서 $F$는 체이고 $K/F$는 원소
$\alpha_1, \ldots, \alpha_n \in K$에 의해 생성되는 유한확대이다.
% P01-E0250: frozen source line 1270 omits the finite verb.
$K_0 = F$로 놓고
$$
K_i = F(\alpha_1, \ldots, \alpha_i)
$$
로 두면 유한확대들의 탑
$K = K_n / K_{n - 1} / \ldots / K_0 = F$를 얻는다.
$P_i$를 $\alpha_i$의 $K_{i - 1}$ 위 최소다항식이라 하자.
마지막으로 $F$의 대수적 폐포 $\overline{F}$를 하나 고정한다.
\end{situation}

\noindent
$F$, $K$, $\alpha_i$, $\overline{F}$는 Situation
\ref{fields-situation-finitely-generated}에서와 같다고 하자.
$\varphi : K \to \overline{F}$가 $F$-확대의 사상이라고 가정하자.
그러면 사상 $\varphi_i : K_i \to \overline{F}$를 얻는다.
$P_i \in K_{i - 1}[x]$의 $\varphi_{i - 1}$에 의한 상은 다항식
$P_i^{\varphi_{i - 1}} \in \overline{F}[x]$이고,
$\varphi(\alpha_i)$는 이 다항식의 근이다.

\begin{lemma}
\label{fields-lemma-count-embeddings}
Situation \ref{fields-situation-finitely-generated}에서 대응
$$
\Mor_F(K, \overline{F})
\longrightarrow
\{(\beta_1, \ldots, \beta_n)\text{ as below}\},
\quad
\varphi \longmapsto (\varphi(\alpha_1), \ldots, \varphi(\alpha_n))
$$
은 전단사이다. 여기서 오른쪽은 다음 성질을 갖는
$\overline{F}$의 원소들의 $n$-튜플
$(\beta_1, \ldots, \beta_n)$의 집합이다.
\begin{enumerate}
\item $\beta_1 \in \overline{F}$는 $P_1$의 근이다.
$\varphi_1 : K_1 \to \overline{F}$를 $\alpha_1$을 $\beta_1$로 보내는
$F$ 위의 준동형이라 하자.
\item $\beta_2 \in \overline{F}$는 $P_2^{\varphi_1}$의 근이다.
$\varphi_2 : K_2 \to \overline{F}$를 $\varphi_1$을 연장하고
$\alpha_2$를 $\beta_2$로 보내는 준동형이라 하자.
\item 이와 같이 계속하여 마지막에는
\item $\beta_n \in \overline{F}$가
$P_n^{\varphi_{n - 1}}$의 근이다.
\end{enumerate}
각 단계에서 준동형 $\varphi_i$가 존재하고 유일한 것은
$K_i = K_{i - 1}[x]/(P_i)$이고 $\beta_i$가
$P_i^{\varphi_{i - 1}}$의 근이기 때문이다.
% P01-E0251: frozen source lines 1313 and 1320 misspell “homomorphism”.
\end{lemma}

\begin{proof}
왼쪽에서 오른쪽으로 가는 사상은 보조정리 앞에서 설명하였다.
$(\beta_1, \ldots, \beta_n)$을 오른쪽의 원소라 하자.
그러면 $\varphi : K = K_n \to \overline{F}$를
$\varphi_{n - 1}$을 연장하고 $\alpha_n$을 $\beta_n$으로 보내는
유일한 준동형으로 둔다. 유일성으로부터 두 구성이 서로 역임을 알 수 있다.
\end{proof}

\begin{lemma}
\label{fields-lemma-count-embeddings-explicitly}
Situation \ref{fields-situation-finitely-generated}에서
$|\Mor_F(K, \overline{F})| = \prod_{i = 1}^n \deg_s(P_i)$이다.
\end{lemma}

\begin{proof}
이는 Lemma \ref{fields-lemma-count-embeddings}에서 즉시 따른다.
여기서 암묵적으로 사용하는 핵심 사실은 Definition
\ref{fields-definition-separable-degree} 바로 앞에서 확인한
기약다항식의 분리차수가 잘 정의된다는 것이다.
\end{proof}

\noindent
이제 위의 결과를 사용하여 분리 가능 체확대를 특징지을 수 있다.

\begin{lemma}
\label{fields-lemma-separably-generated-separable}
가정과 표기는 Situation \ref{fields-situation-finitely-generated}에서와 같다.
각 $P_i$가 분리 가능, 즉 각 $\alpha_i$가
$K_{i - 1}$ 위에서 분리 가능하면
$$
|\Mor_F(K, \overline{F})| = [K : F]
$$
이고 체확대 $K/F$는 분리 가능하다. 어느 $\alpha_i$가
$K_{i - 1}$ 위에서 분리 가능하지 않으면
$|\Mor_F(K, \overline{F})| < [K : F]$이다.
\end{lemma}

\begin{proof}
$\alpha_i$가 $K_{i - 1}$ 위에서 분리 가능하면
$\deg_s(P_i) = \deg(P_i) = [K_i : K_{i - 1}]$이다
(마지막 등식은 Lemma \ref{fields-lemma-degree-minimal-polynomial}에 따른다).
곱셈성(Lemma \ref{fields-lemma-multiplicativity-degrees})에 의해
$$
[K : F] = \prod [K_i : K_{i - 1}] = \prod \deg(P_i) =
\prod \deg_s(P_i) = |\Mor_F(K, \overline{F})|
$$
이다. 마지막 등식은 Lemma
\ref{fields-lemma-count-embeddings-explicitly}에 따른다.
똑같은 논증으로, 어느 $\alpha_i$가 $K_{i - 1}$ 위에서
분리 가능하지 않으면
$|\Mor_F(K, \overline{F})| < [K : F]$라는 진부등식을 얻는다.

\medskip\noindent
마지막으로 다시 각 $\alpha_i$가 $K_{i - 1}$ 위에서
분리 가능하다고 가정하자. $K/F$가 분리 가능함을 보이겠다.
임의의 $\gamma = \gamma_1 \in K$를 택하자. 그러면
$K = F(\gamma_1, \ldots, \gamma_m)$이 되도록 원소
$\gamma_2, \ldots, \gamma_m$을 더 택할 수 있다(예를 들어
$\gamma_2 = \alpha_1, \ldots, \gamma_{n + 1} = \alpha_n$으로
택할 수 있다). 이미 위에서 증명한 보조정리의 마지막 부분에 따르면,
$\gamma$가 $F$ 위에서 분리 가능하지 않을 경우
$|\Mor_F(K, \overline{F})| < [K : F]$라는 진부등식을 얻는다.
이는 역시 이미 위에서 증명한 보조정리의 첫 부분과 모순이다.
% P01-E0252: frozen source lines 1379--1380 say “already prove above”.
\end{proof}

\begin{lemma}
\label{fields-lemma-separable-equality}
$K/F$를 체의 유한확대라 하고 $\overline{F}$를 $F$의
대수적 폐포라 하자. 그러면
$$
|\Mor_F(K, \overline{F})| \leq [K : F]
$$
이고, 등식이 성립할 필요충분조건은 $K$가 $F$ 위에서
분리 가능한 것이다.
\end{lemma}

\begin{proof}
이는 Lemma \ref{fields-lemma-separably-generated-separable}의 따름정리이다.
실제로 $K/F$가 유한이므로 $F$ 위에서 $K$를 생성하는 유한개의 원소
$\alpha_1, \ldots, \alpha_n \in K$를 찾을 수 있다(예를 들어
$\alpha_i$들을 $F$ 위에서 $K$의 기저로 택할 수 있다).
$K/F$가 분리 가능하면 Lemma \ref{fields-lemma-separable-goes-up}에 의해
각 $\alpha_i$가 $F(\alpha_1, \ldots, \alpha_{i - 1})$
위에서 분리 가능하고, Lemma
\ref{fields-lemma-separably-generated-separable}에 의해 등식을 얻는다.
한편 등식이 성립하면 $\alpha_1, \ldots, \alpha_n$을 어떻게
택하더라도 Lemma \ref{fields-lemma-separably-generated-separable}에 의해
$\alpha_1$은 $F$ 위에서 분리 가능하다. 이 열을 $K$의 임의의
원소로 시작할 수 있으므로 $K$는 $F$ 위에서 분리 가능하다.
\end{proof}

\begin{lemma}
\label{fields-lemma-separable-permanence}
$E/k$와 $F/E$를 분리 가능 대수적 체확대라 하자. 그러면 $F/k$는
분리 가능 체확대이다.
\end{lemma}

\begin{proof}
$\alpha \in F$를 택하자. 그러면 $\alpha$는 $E$ 위에서 분리 가능하고
대수적이다. $P = x^d + \sum_{i < d} a_i x^i$를 $\alpha$의
$E$ 위 최소다항식이라 하자. 각 $a_i$는 $k$ 위에서 분리 가능하고
대수적이다. 체들의 탑
$$
k \subset k(a_0) \subset k(a_0, a_1) \subset \ldots \subset
k(a_0, \ldots, a_{d - 1}) \subset k(a_0, \ldots, a_{d - 1}, \alpha)
$$
을 생각하자. $a_i$가 $k$ 위에서 분리 가능하고 대수적이므로
Lemma \ref{fields-lemma-separable-goes-up}에 의해
$k(a_0, \ldots, a_{i - 1})$ 위에서도 분리 가능하고 대수적이다.
끝으로 $\alpha$는 $k(a_0, \ldots, a_{d - 1})$ 위에서 분리 가능하고
대수적이다. 실제로 이는 $P$의 근이며, 이 다항식은
기약이고(더 클 수도 있는 체 $E$ 위에서도 기약이므로)
분리 가능하다($E$ 위에서 분리 가능하므로).
따라서 Lemma \ref{fields-lemma-separably-generated-separable}에 의해
$k(a_0, \ldots, a_{d - 1}, \alpha)$는 $k$ 위에서 분리 가능하며,
원하는 대로 $\alpha$가 $k$ 위에서 분리 가능하다는 결론을 얻는다.
\end{proof}

\begin{lemma}
\label{fields-lemma-separable-elements}
$E/k$를 체확대라 하자. 그러면 $E$에서 $k$ 위에 분리 가능한 원소들은
$E/k$의 부분확대를 이룬다.
\end{lemma}

\begin{proof}
$\alpha, \beta \in E$가 $k$ 위에서 분리 가능하다고 하자.
Lemma \ref{fields-lemma-separable-goes-up}에 의해 $\beta$는
$k(\alpha)$ 위에서 분리 가능하다. Lemma
\ref{fields-lemma-separably-generated-separable}를 $n = 2$,
$\alpha_1 = \alpha$, $\alpha_2 = \beta$로 적용하면
$k(\alpha, \beta)$가 $k$ 위에서 분리 가능함을 알 수 있다.
\end{proof}





\section{지표의 선형독립성}
\label{fields-section-independence-characters}

\noindent
필요한 명제는 다음과 같다.

\begin{lemma}
\label{fields-lemma-independence-characters}
$L$을 체라 하고 $G$를 모노이드, 예를 들어 군이라 하자.
$\chi_1, \ldots, \chi_n : G \to L$을 서로 다른 모노이드 준동형이라
하자. 여기서 $L$은 곱셈에 의해 모노이드로 여긴다.
그러면 $\chi_1, \ldots, \chi_n$은 $L$ 위에서 선형독립이다.
즉 $\lambda_1, \ldots, \lambda_n \in L$이 모두 영은 아니면,
어떤 $g \in G$에 대하여 $\sum \lambda_i\chi_i(g) \not = 0$이다.
\end{lemma}

\begin{proof}
$n = 1$일 때는 $e \in G$가 중립원(항등원)이면
$\chi_1(e) = 1$이므로 참이다. $n > 1$에 대하여 귀납법으로
결과를 증명하자. $\lambda_1, \ldots, \lambda_n \in L$이 모두
영은 아니라고 하자. 어떤 항에 대하여 $\lambda_i = 0$이면
$n$에 관한 귀납법으로 결론을 얻는다. 어떤 $g \in G$에 대하여
$\sum \lambda_i\chi_i(g) \not = 0$임을 보이려는 것이므로
$-\lambda_n$으로 나누어 $\lambda_n = -1$이라 가정할 수 있다.
문제가 생길 수 있는 유일한 경우는
$$
\chi_n(g) = \sum\nolimits_{i = 1, \ldots, n - 1} \lambda_i\chi_i(g)
$$
가 모든 $g \in G$에 대하여 성립하는 때이다. $h \in G$를
고정하자. 그러면 또한
\begin{align*}
\chi_n(h)\chi_n(g) & = \chi_n(hg) \\
& = \sum\nolimits_{i = 1, \ldots, n - 1} \lambda_i\chi_i(hg) \\
& = \sum\nolimits_{i = 1, \ldots, n - 1} \lambda_i\chi_i(h) \chi_i(g)
\end{align*}
를 얻는다. 앞의 관계에 $\chi_n(h)$를 곱하고 빼면
$$
0 = \sum\nolimits_{i = 1, \ldots, n - 1}
\lambda_i (\chi_n(h) - \chi_i(h)) \chi_i(g)
$$
를 모든 $g \in G$에 대하여 얻는다. $\lambda_i \not = 0$이므로
귀납법에 의해 모든 $i$에 대하여
$\chi_n(h) = \chi_i(h)$라는 결론을 얻는다. 위에서 $h$는
임의로 택했으므로 $i \leq n - 1$에 대하여
$\chi_i = \chi_n$이다. 이는 지표 $\chi_i$들이 서로 다르다는
가정과 모순이다.
\end{proof}

\begin{lemma}
\label{fields-lemma-sums-of-powers}
$L$을 체라 하자. $n \geq 1$이고
$\alpha_1, \ldots, \alpha_n \in L$이 $L$의 서로 다른
원소들이라고 하자. 그러면 어떤 $e \geq 0$에 대하여
$\sum_{i = 1, \ldots, n} \alpha_i^e \not = 0$이다.
\end{lemma}

\begin{proof}
지표의 선형독립성(Lemma \ref{fields-lemma-independence-characters})을
모노이드 준동형 $\mathbf{Z}_{\geq 0} \to L$,
$e \mapsto \alpha_i^e$에 적용한다.
\end{proof}

\begin{lemma}
\label{fields-lemma-independence-embeddings}
$K/F$와 $L/F$를 체확대라 하자.
$\sigma_1, \ldots, \sigma_n : K \to L$을 서로 다른
$F$-확대의 사상이라 하자. 그러면
$\sigma_1, \ldots, \sigma_n$은 $L$ 위에서 선형독립이다.
즉 $\lambda_1, \ldots, \lambda_n \in L$이 모두 영은 아니면,
어떤 $\alpha \in K$에 대하여
$\sum \lambda_i\sigma_i(\alpha) \not = 0$이다.
\end{lemma}

\begin{proof}
Lemma \ref{fields-lemma-independence-characters}를 단위원군에 대한
$\sigma_i$들의 제한에 적용한다.
\end{proof}

\begin{lemma}
\label{fields-lemma-finite-separable-tensor-alg-closed}
$K/F$와 $L/F$를 체확대라 하고, $K/F$는 유한 분리 가능 확대이며
$L$은 대수적으로 닫혀 있다고 하자. 그러면 사상
$$
K \otimes_F L
\longrightarrow
\prod\nolimits_{\sigma \in \Hom_F(K, L)} L,\quad
\alpha \otimes \beta \mapsto (\sigma(\alpha)\beta)_\sigma
$$
은 $L$-대수의 동형사상이다.
\end{lemma}

\begin{proof}
$\alpha_1, \ldots, \alpha_n$을 $F$ 위의 벡터공간으로서
$K$의 기저로 택하자. Lemma \ref{fields-lemma-separable-equality}와
생략된 짧은 논증에 의해 집합 $\Hom_F(K, L)$은 $n$개의 원소,
이를테면 $\sigma_1, \ldots, \sigma_n$을 갖는다. 특히 양변은
$L$ 위의 벡터공간으로서 같은 차원 $n$을 갖는다.
따라서 이 사상이 동형사상이 아니면 핵을 갖는다. 다시 말해,
$\mu_j \in L$, $j = 1, \ldots, n$ 중 모두가 영은 아니면서
$\sum \alpha_j \otimes \mu_j$가 핵에 속하는 것들이 존재한다.
즉 모든 $i$에 대하여 $\sum \sigma_i(\alpha_j)\mu_j = 0$이다.
이는 성분이 $\sigma_i(\alpha_j)$인 $n \times n$ 행렬이
가역이 아님을 뜻한다. 따라서 모두가 영은 아닌
$\lambda_1, \ldots, \lambda_n \in L$을 찾아 모든 $j$에 대하여
$\sum \lambda_i\sigma_i(\alpha_j) = 0$이 되게 할 수 있다.
이제 임의의 원소 $\alpha \in K$는 $\beta_j \in F$를 사용하여
$\alpha = \sum \beta_j \alpha_j$로 쓸 수 있고, 그러면
$$
\sum \lambda_i\sigma_i(\alpha) =
\sum \lambda_i\sigma_i(\sum \beta_j \alpha_j) =
\sum \beta_j \sum \lambda_i\sigma_i(\alpha_j) = 0
$$
을 얻는다. 이는 Lemma \ref{fields-lemma-independence-embeddings}와 모순이다.
\end{proof}






\section{순수 비분리 확대}
\label{fields-section-purely-inseparable}

\noindent
순수 비분리 확대는 앞 절에서 정의한 분리 가능 확대와 반대되는
것이다. 이러한 확대는 양의 표수에서만 나타난다.

\begin{definition}
\label{fields-definition-purely-inseparable}
$F$를 표수가 $p > 0$인 체라 하고 $K/F$를 확대라 하자.
\begin{enumerate}
\item 어떤 $p$의 거듭제곱 $q$에 대하여 $\alpha^q \in F$이면
원소 $\alpha \in K$가 $F$ 위에서 {\it 순수 비분리}라고 한다.
\item 확대 $K/F$가 {\it 순수 비분리}라는 것은 $K$의 모든 원소가
$F$ 위에서 순수 비분리라는 것과 동치이다.
\end{enumerate}
체확대 $L/M$이 주어졌다고 하자(여기서는 $M$의 표수에 아무
조건도 없다). $L = M$이거나 $M$의 표수가 소수 $p$이고
$L/M$이 위에서 정의한 의미로 순수 비분리이면 이 확대를
{\it 순수 비분리}라고 한다.
\end{definition}

\noindent
순수 비분리 확대는 반드시 대수적임에 유의하자.
$F$를 표수가 $p > 0$인 체라 하자. 아직 $p$제곱근을 갖지 않는
원소 $t \in F$에 그 근을 첨가하면 순수 비분리 확대의 한 예를
얻는다. 실제로 아래 보조정리는 $P = x^p - t$가 기약임을 보이며,
따라서
$$
K = F[x]/(P) = F[t^{1/p}]
$$
는 체이다. 또한 $K$는 $F$ 위에서 순수 비분리이다. 왜냐하면
$K$의 모든 원소
$$
a_0 + a_1t^{1/p} + \ldots + a_{p - 1}t^{(p - 1)/p},\quad a_i \in F
$$
의 $p$제곱은
$$
(a_0 + a_1t^{1/p} + \ldots + a_{p - 1}t^{(p - 1)/p})^p =
a_0^p + a_1^p t + \ldots + a_{p - 1}^pt^{p - 1} \in F
$$
와 같기 때문이다. 이러한 상황은 $\mathbf{F}_p$ 위의
유리함수체 $\mathbf{F}_p(t)$에서 일어난다.

\begin{lemma}
\label{fields-lemma-take-pth-root}
$p$를 소수라 하고 $F$를 표수가 $p$인 체라 하자.
$t \in F$가 $F$ 안에서 $p$제곱근을 갖지 않는 원소라고 하자.
그러면 다항식 $x^p - t$는 $F$ 위에서 기약이다.
\end{lemma}

\begin{proof}
이를 보이기 위해 인수분해 $x^p - t = f g$가 있다고 가정하자.
미분하면 $f' g + f g' = 0$을 얻는다. $f$와 $g$의 차수는
$p$보다 작으므로 $f' = 0$도 $g' = 0$도 아님에 유의하자.
더욱이 $\deg(f') < \deg(f)$이고 $\deg(g') < \deg(g)$이다.
따라서 $f$와 $g$는 공통인수를 갖는다. 그러므로 $x^p - t$가
가약이면 어떤 기약다항식 $f$, $c \in F^*$, $n > 1$에 대하여
$x^p - t = c f^n$의 꼴이다. $p$가 소수이므로 이는 $n = p$이고
$f$가 일차임을 뜻하며, 따라서 $x^p - t$가 $F$ 안에 근을
가져야 한다. 이는 모순이다.
\end{proof}

\noindent
표수 $p$에서 $p$제곱근을 취하는 일이 매우 중요한 연산임을
보게 될 것이다.

\begin{lemma}
\label{fields-lemma-purely-inseparable-permanence}
$E/k$와 $F/E$를 순수 비분리 체확대라 하자. 그러면 $F/k$는
순수 비분리 체확대이다.
\end{lemma}

\begin{proof}
$k$의 표수가 $p$라고 하자. $\alpha \in F$를 택하자.
어떤 $p$의 거듭제곱 $q$에 대하여 $\alpha^q \in E$이다.
그러면 어떤 $p$의 거듭제곱 $q'$에 대하여
$(\alpha^q)^{q'} \in k$이다. 따라서 $\alpha^{qq'} \in k$이다.
\end{proof}

\begin{lemma}
\label{fields-lemma-purely-inseparable-elements}
$E/k$를 체확대라 하자. 그러면 $E$에서 $k$ 위에 순수 비분리인
원소들은 $E/k$의 부분확대를 이룬다.
\end{lemma}

\begin{proof}
$p$를 $k$의 표수라 하자. $\alpha, \beta \in E$가 $k$ 위에서
순수 비분리라고 하자. 어떤 $p$의 거듭제곱 $q, q'$에 대하여
$\alpha^q \in k$이고 $\beta^{q'} \in k$이라고 하자.
$q''$이 $p$의 거듭제곱이면
$(\alpha + \beta)^{q''} = \alpha^{q''} + \beta^{q''}$이다.
따라서 $q'' \geq q, q'$이면 $\alpha + \beta$가 $k$ 위에서
순수 비분리이다. $\alpha$와 $\beta$의 차, 곱, 몫도 마찬가지이다.
\end{proof}

\begin{lemma}
\label{fields-lemma-finite-purely-inseparable}
$E/F$를 표수가 $p > 0$인 유한 순수 비분리 체확대라 하자.
그러면 원소들의 열 $\alpha_1, \ldots, \alpha_n \in E$가 존재하여
체들의 탑
$$
E = F(\alpha_1, \ldots, \alpha_n) \supset
F(\alpha_1, \ldots, \alpha_{n - 1}) \supset
\ldots
\supset F(\alpha_1) \supset F
$$
을 얻는다. 각 중간 확대의 차수는 $p$이고 $p$제곱근 하나를
첨가하여 얻어진다. 구체적으로
$\alpha_i^p \in F(\alpha_1, \ldots, \alpha_{i - 1})$는
$F(\alpha_1, \ldots, \alpha_{i - 1})$ 안에서 $p$제곱근을
갖지 않는 원소이다. 여기서 $i = 1, \ldots, n$이다.
\end{lemma}

\begin{proof}
$E/F$의 차수에 관한 귀납법을 사용한다. 확대의 차수가 $1$이면
결과는 자명하다($n = 0$으로 둔다). 그렇지 않으면
$\alpha \in E$, $\alpha \not \in F$를 택하자. 어떤 $r > 0$에
대하여 $\alpha^{p^r} \in F$라고 하자. $r$을 최소로 택하고
$\alpha$를 $\alpha^{p^{r - 1}}$로 대체한다. 그러면
$\alpha \not \in F$이지만 $\alpha^p \in F$이다.
이제 $t = \alpha^p$는 $F$ 안의 $p$제곱이 아니다(그렇다면
$\alpha \in F$가 되기 때문이다. Lemma
\ref{fields-lemma-nr-roots-unchanged} 또는 그 증명을 보라).
따라서 $F \subset F(\alpha)$는 차수가 $p$인 부분확대이다
(Lemma \ref{fields-lemma-take-pth-root}). 귀납법에 의해 보조정리의
결론을 만족하면서 $E/F(\alpha)$를 생성하는
$\alpha_1, \ldots, \alpha_n \in E$를 찾을 수 있다.
열 $\alpha, \alpha_1, \ldots, \alpha_n$이 확대 $E/F$에 대하여
원하는 역할을 한다.
\end{proof}

\begin{lemma}
\label{fields-lemma-separable-first}
\begin{slogan}
모든 대수적 체확대는 유일하게 분리 가능 체확대에 순수 비분리
체확대가 뒤따르는 꼴이다.
\end{slogan}
$E/F$를 대수적 체확대라 하자. $E_{sep}/F$가 분리 가능하고
$E/E_{sep}$가 순수 비분리인 유일한 부분확대
$E/E_{sep}/F$가 존재한다.
\end{lemma}

\begin{proof}
표수가 영이면 $E_{sep} = E$로 둔다. 표수가 $p > 0$이라고
가정하자. $E_{sep}$를 $E$에서 $F$ 위에 분리 가능한 원소들의
집합이라 하자. Lemma \ref{fields-lemma-separable-elements}에 의해
이는 부분확대이고, 물론 $E_{sep}$는 $F$ 위에서 분리 가능하다.
$E$의 원소 $\alpha$가 주어지면 어떤 $p$의 거듭제곱 $q$에
대하여 $\alpha^q$는 $F$ 위에서 분리 가능하다.
구체적으로 $q$는 $\alpha$의 최소다항식이 $P(x^q)$의 꼴이고
$P$가 분리 가능한 대수적 대상이 되게 하는 $p$의 거듭제곱이다.
% P01-E1267: frozen source calls the polynomial “separable algebraic”.
Lemma \ref{fields-lemma-irreducible-polynomials}를 보라. 따라서
$E/E_{sep}$는 순수 비분리이다. 유일성은 자명하다.
\end{proof}

\begin{definition}
\label{fields-definition-insep-degree}
$E/F$를 대수적 체확대라 하고 $E_{sep}$를 Lemma
\ref{fields-lemma-separable-first}에서 얻은 부분확대라 하자.
\begin{enumerate}
\item 정수 $[E_{sep} : F]$를 이 확대의 {\it 분리차수}라 한다.
표기는 $[E : F]_s$이다.
\item 정수 $[E : E_{sep}]$를 이 확대의 {\it 비분리차수}, 또는
{\it 비분리도}라 한다. 표기는 $[E : F]_i$이다.
\end{enumerate}
% P01-E1268: frozen source calls possibly infinite degrees integers.
\end{definition}

\noindent
물론 표수 $0$에서는 $[E : F] = [E : F]_s$이고
$[E : F]_i = 1$이다. 곱셈성(Lemma
\ref{fields-lemma-multiplicativity-degrees})에 의해
$$
[E : F] = [E : F]_s [E : F]_i
$$
이다. 이 등식은 이들 차수 가운데 일부가 무한인 경우에도 성립한다.
실제로 분리차수와 비분리차수 역시 곱셈적이다(Lemma
\ref{fields-lemma-multiplicativity-all-degrees}를 보라).

\begin{lemma}
\label{fields-lemma-separable-degree}
$K/F$를 유한확대라 하고 $\overline{F}$를 $F$의 대수적 폐포라
하자. 그러면 $[K : F]_s = |\Mor_F(K, \overline{F})|$이다.
\end{lemma}

\begin{proof}
먼저 $K/F$가 순수 비분리인 경우를 증명한다. 이 경우에는 유일한
사상 $K \to \overline{F}$가 있다고 주장한다. 이는 Lemma
\ref{fields-lemma-finite-purely-inseparable}에서와 같은 원소들의 열
$\alpha_1, \ldots, \alpha_n \in K$를 택함으로써 알 수 있다.
$\alpha_i$의 $F(\alpha_1, \ldots, \alpha_{i - 1})$ 위
기약다항식은 $x^p - \alpha_i^p$이다. Lemma
\ref{fields-lemma-count-embeddings-explicitly}를 적용하면
$|\Mor_F(K, \overline{F})| = 1$임을 알 수 있다. 한편 이 경우
$[K : F]_s = 1$이므로 등식이 성립한다.

\medskip\noindent
일반적인 유한확대 $K/F$로 돌아가자. 이 경우 Lemma
\ref{fields-lemma-separable-first}에서와 같이 $F \subset K_s \subset K$를
택한다. Lemma \ref{fields-lemma-separable-equality}에 의해
$|\Mor_F(K_s, \overline{F})| = [K_s : F] = [K : F]_s$이다.
한편 앞 단락의 결과에 의해 모든 체 사상
$\sigma' : K_s \to \overline{F}$는 유일한 체 사상
$\sigma : K \to \overline{F}$로 연장된다. 다시 말해
$|\Mor_F(K, \overline{F})| = |\Mor_F(K_s, \overline{F})|$이고,
증명이 끝난다.
\end{proof}

\begin{lemma}[곱셈성]
\label{fields-lemma-multiplicativity-all-degrees}
대수적 체확대들의 탑 $K/E/F$가 주어졌다고 하자. 그러면
$$
[K : F]_s = [K : E]_s [E : F]_s
\quad\text{and}\quad
[K : F]_i = [K : E]_i [E : F]_i
$$
이다.
\end{lemma}

\begin{proof}
먼저 $K$가 $F$ 위에서 유한인 경우를 증명한다.
보통 차수의 곱셈성(Lemma \ref{fields-lemma-multiplicativity-degrees})이
있으므로 두 공식 가운데 하나만 증명하면 충분하다. Lemma
\ref{fields-lemma-separable-degree}에 의해
$[K : F]_s = |\Mor_F(K, \overline{F})|$이다. 같은 보조정리에 의해
임의의 $\sigma \in \Mor_F(E, \overline{F})$가 주어졌을 때,
$\sigma$를 사상 $\tau : K \to \overline{F}$로 연장하는 방법의
수는 $[K : E]_s$이다. 실제로
$E \cong \sigma(E) \subset \overline{F}$를 통하여
$\overline{F}$를 $E$의 대수적 폐포로 볼 수 있다.
$\sigma$를 택하는 방법이
$[E : F]_s = |\Mor_F(E, \overline{F})|$개라는 사실과 합치면
결과를 얻는다.

\medskip\noindent
확대가 무한인 경우의 증명은 생략한다.
\end{proof}








\section{정규확대}
\label{fields-section-normal}

\noindent
$P \in F[x]$를 체 $F$ 위의 비상수다항식이라 하자.
어떤 $c \in F^*$, $n \geq 1$, $\alpha_1, \ldots, \alpha_n \in F$에
대하여
$$
P = c(x - \alpha_1) \ldots (x - \alpha_n)
$$
가 $F[x]$에서 성립하면 $P$가 {\it $F$ 위에서 일차인수들로
완전히 분해된다} 또는 {\it $F$ 위에서 완전히 분해된다}고 한다.
정규확대는 다음과 같이 정의한다.

\begin{definition}
\label{fields-definition-normal}
$E/F$를 대수적 체확대라 하자. 모든 $\alpha \in E$에 대하여
$\alpha$의 $F$ 위 최소다항식 $P$가 $E$ 위에서 일차인수들로
완전히 분해되면 $E$가 $F$ 위에서 {\it 정규}라고 한다.
\end{definition}

\noindent
분리 가능 확대의 경우와 마찬가지로, 이 개념의 기본 성질을
확립하려면 약간의 작업이 필요하다.

\begin{lemma}
\label{fields-lemma-normal-goes-up}
$K/E/F$를 대수적 체확대들의 탑이라 하자. $K$가 $F$ 위에서
정규이면 $K$는 $E$ 위에서도 정규이다.
\end{lemma}

\begin{proof}
$\alpha \in K$라 하자. $P$를 $\alpha$의 $F$ 위 최소다항식,
$Q$를 $\alpha$의 $E$ 위 최소다항식이라 하자.
그러면 $Q$는 다항식환 $E[x]$에서 $P$를 나누므로
$P = QR$이라 쓸 수 있다. 따라서 $P$가 $K$ 위에서 완전히
분해되면 $Q$도 그러하다.
\end{proof}

\begin{lemma}
\label{fields-lemma-intersect-normal}
$F$를 체라 하고 $M/F$를 대수적 확대라 하자.
$M/E_i/F$, $i \in I$를 $E_i/F$가 정규인 부분확대들이라 하자.
그러면 $\bigcap E_i$는 $F$ 위에서 정규이다.
\end{lemma}

\begin{proof}
정의에서 바로 따른다.
\end{proof}

\begin{lemma}
\label{fields-lemma-separable-first-normal}
$E/F$를 정규 대수적 체확대라 하자. 그러면 Lemma
\ref{fields-lemma-separable-first}의 부분확대 $E/E_{sep}/F$는 정규이다.
\end{lemma}

\begin{proof}
표수가 영이면 $E_{sep} = E$이고 결과는 자명하다.
표수가 $p > 0$이면 $E_{sep}$는 $E$에서 $F$ 위에 분리 가능한
원소들의 집합이다. 이제 $\alpha \in E_{sep}$의 최소다항식을
$P$라 하고
$P = c(x - \alpha)(x - \alpha_2) \ldots (x - \alpha_d)$로 쓰자.
여기서 $\alpha_2, \ldots, \alpha_d \in E$이다.
$P$가 분리 가능 다항식이고 $\alpha_i$가 $P$의 근이므로,
원하는 대로 $\alpha_i \in E_{sep}$라는 결론을 얻는다.
\end{proof}

\begin{lemma}
\label{fields-lemma-characterize-normal}
$E/F$를 체의 대수적 확대라 하고 $\overline{F}$를 $F$의
대수적 폐포라 하자. 다음은 동치이다.
\begin{enumerate}
\item $E$는 $F$ 위에서 정규이다.
\item 모든 쌍 $\sigma, \sigma' \in \Mor_F(E, \overline{F})$에 대하여
$\sigma(E) = \sigma'(E)$이다.
\end{enumerate}
\end{lemma}

\begin{proof}
$\mathcal{P}$를 $E$의 모든 원소들의 $F$ 위 최소다항식 전체의
집합이라 하자. 다음과 같이 놓는다.
$$
T =
\{\beta \in \overline{F} \mid P(\beta) = 0\text{ for some }P \in \mathcal{P}\}
$$
$E$가 $F$ 위에서 정규이면 모든
$\sigma \in \Mor_F(E, \overline{F})$에 대하여
$\sigma(E) = T$임은 자명하다. 따라서 (1)은 (2)를 함의한다.

\medskip\noindent
거꾸로 (2)를 가정하고 $\beta \in T$를 택하자.
최소다항식 $P \in \mathcal{P}$가 $\beta$를 영으로 만드는
해당 원소 $\alpha \in E$를 찾을 수 있다.
$F(\alpha) = F[x]/(P)$이므로 $\alpha$를 $\beta$로 보내는 원소
$\sigma_0 \in \Mor_F(F(\alpha), \overline{F})$를 찾을 수 있다.
Lemma \ref{fields-lemma-map-into-algebraic-closure}에 의해
$\sigma_0$를 $\sigma \in \Mor_F(E, \overline{F})$로 연장할 수 있다.
따라서 $\beta$는 모든 매입 $\sigma : E \to \overline{F}$의
공통 상에 속한다. 그러므로 임의의 $\sigma$에 대하여
$\sigma(E) = T$이다. 하나의 $\sigma$를 고정하자.
이제 $P \in \mathcal{P}$라 하자. Lemma
\ref{fields-lemma-algebraically-closed}에 의해 어떤 $n$과
$\beta_i \in \overline{F}$에 대하여
$$
P = (x - \beta_1) \ldots (x - \beta_n)
$$
으로 쓸 수 있다. $\beta_i \in T$임에 유의하자.
따라서 어떤 $\alpha_i \in E$에 대하여
$\beta_i = \sigma(\alpha_i)$이다. 그러므로
$P = (x - \alpha_1) \ldots (x - \alpha_n)$은 $E$ 위에서
완전히 분해된다. 이로써 증명이 끝난다.
\end{proof}

\begin{lemma}
\label{fields-lemma-normally-generated}
$E/F$를 체의 대수적 확대라 하자. $E$가 $F$ 위에서
$\alpha_i \in E$, $i \in I$에 의해 생성되고, 각 $i$에 대하여
$\alpha_i$의 $F$ 위 최소다항식이 $E$에서 완전히 분해된다고 하자.
그러면 $E/F$는 정규이다.
\end{lemma}

\begin{proof}
$P_i$를 $\alpha_i$의 $F$ 위 최소다항식이라 하자.
$\alpha_i = \alpha_{i, 1}, \alpha_{i, 2}, \ldots, \alpha_{i, d_i}$를
$P_i$의 $E$ 위 근들이라 하자. $F$ 위의 두 매입
$\sigma, \sigma' : E \to \overline{F}$가 주어지면
$$
\{\sigma(\alpha_{i, 1}), \ldots, \sigma(\alpha_{i, d_i})\} =
\{\sigma'(\alpha_{i, 1}), \ldots, \sigma'(\alpha_{i, d_i})\}
$$
이다. 양변이 모두 $\overline{F}$에서 $P_i$의 근들의 집합과
같기 때문이다. 원소들 $\alpha_{i, j}$가 $F$ 위에서 $E$를
생성하므로 $\sigma(E) = \sigma'(E)$이다. 따라서 Lemma
\ref{fields-lemma-characterize-normal}에 의해 $E/F$는 정규이다.
\end{proof}

\begin{lemma}
\label{fields-lemma-lift-maps}
$L/M/K$를 대수적 확대들의 탑이라 하자.
\begin{enumerate}
\item $M/K$가 정규이면 $L/K$의 모든 자기동형사상 $\tau$는
자기동형사상 $\tau|_M : M \to M$을 유도한다.
\item $L/K$가 정규이면 모든 $K$-대수 사상
$\sigma : M \to L$은 $L$의 자기동형사상으로 연장된다.
\end{enumerate}
\end{lemma}

\begin{proof}
$L$의 대수적 폐포 $\overline{L}$을 택한다
(Theorem \ref{fields-theorem-existence-algebraic-closure}).

\medskip\noindent
$\tau$가 (1)에서와 같다고 하자. Lemma
\ref{fields-lemma-characterize-normal}에 의해 $\overline{L}$의
부분체들로서 $\tau(M) = M$이고, 따라서
$\tau|_M : M \to M$은 자기동형사상이다.

\medskip\noindent
$\sigma : M \to L$이 (2)에서와 같다고 하자. Lemma
\ref{fields-lemma-map-into-algebraic-closure}에 의해 $\sigma$를 사상
$\tau : L \to \overline{L}$로 연장할 수 있다. 즉
$$
\xymatrix{
L \ar[r]_\tau & \overline{L} \\
M \ar[u] \ar[ru]_\sigma & K \ar[l] \ar[u]
}
$$
가 가환이다. Lemma \ref{fields-lemma-characterize-normal}에 의해
$\tau(L) = L$이다. 따라서 $\tau : L \to L$은 $\sigma$를
연장하는 자기동형사상이다.
\end{proof}

\begin{definition}
\label{fields-definition-automorphisms}
$E/F$를 체확대라 하자. $\text{Aut}(E/F)$ 또는
$\text{Aut}_F(E)$는 $F$-확대의 범주의 대상으로서 $E$의
자기동형사상군을 나타낸다. $\text{Aut}(E/F)$의 원소들을
{\it $F$ 위의 $E$의 자기동형사상} 또는
{\it $E/F$의 자기동형사상}이라 한다.
\end{definition}

\noindent
자기동형사상을 사용한 정규확대의 특징짓기는 다음과 같다.

\begin{lemma}
\label{fields-lemma-normal-and-automorphisms}
$E/F$를 유한확대라 하자. 그러면
$$
|\text{Aut}(E/F)| \leq [E : F]_s
$$
이고, 등식이 성립할 필요충분조건은 $E$가 $F$ 위에서
정규인 것이다.
\end{lemma}

\begin{proof}
$F$의 대수적 폐포 $\overline{F}$를 택하자.
$[E : F]_s = |\Mor_F(E, \overline{F})|$임을 상기하자.
원소 $\sigma_0 \in \Mor_F(E, \overline{F})$를 택한다.
그러면 사상
$$
\text{Aut}(E/F) \longrightarrow \Mor_F(E, \overline{F}),\quad
\tau \longmapsto \sigma_0 \circ \tau
$$
은 단사이다. 이로써 부등식을 얻는다. 등식이 성립하면 모든
$\sigma \in \Mor_F(E, \overline{F})$는 $\sigma_0$에
자기동형사상을 앞합성하여 얻어진다. 따라서
$\sigma(E) = \sigma_0(E)$이다. 그러므로 Lemma
\ref{fields-lemma-characterize-normal}에 의해 $E$는 $F$ 위에서 정규이다.

\medskip\noindent
거꾸로 $E/F$가 정규라고 가정하자. 그러면 Lemma
\ref{fields-lemma-characterize-normal}에 의해 모든
$\sigma \in \Mor_F(E, \overline{F})$에 대하여
$\sigma(E) = \sigma_0(E)$이다. 따라서
$\tau = \sigma_0^{-1} \circ \sigma$로 놓으면 $F$ 위의 $E$의
자기동형사상을 얻는다. 그러므로 위에 표시한 사상은 전사이다.
\end{proof}

\begin{lemma}
\label{fields-lemma-normal-embeddings-differ-by-aut}
$L/K$를 체의 대수적 정규확대라 하고 $E/K$를 체확대라 하자.
그러면 $L$에서 $E$로 가는 $K$-매입이 존재하지 않거나,
하나의 $\tau : L \to E$가 존재하고 다른 모든 매입은
$\sigma \in \text{Aut}(L/K)$에 대하여 $\tau \circ \sigma$의
꼴이다.
\end{lemma}

\begin{proof}
$\tau$가 주어지면 $L$을 $\tau(L) \subset E$로 대체하고
Lemma \ref{fields-lemma-lift-maps}를 적용한다.
\end{proof}





\section{분해체}
\label{fields-section-splitting-fieds}
% P01-E0253: frozen source label misspells “fields” as “fieds”.

\noindent
다음 보조정리는 정규 체확대를 구성하는 데 유용한 도구이다.

\begin{lemma}
\label{fields-lemma-splitting-field}
$F$를 체라 하고 $P \in F[x]$를 비상수다항식이라 하자.
$P$가 $E$ 위에서 완전히 분해되게 하는 가장 작은 체확대
$E/F$가 존재한다. 더욱이 체확대 $E/F$는 정규이고
(유일하지 않을 수 있는) 동형사상까지 유일하다.
\end{lemma}

\begin{proof}
대수적 폐포 $\overline{F}$를 택하자. 그러면
$P = c (x - \beta_1) \ldots (x - \beta_n)$으로
$\overline{F}[x]$ 안에서 쓸 수 있다. Lemma
\ref{fields-lemma-algebraically-closed}를 보라. $c \in F^*$임에 유의하자.
$E = F(\beta_1, \ldots, \beta_n)$으로 놓는다. 그러면 $P$가
$E$ 위에서 완전히 분해된다는 조건 아래 $E$가 최소임은 자명하다.

\medskip\noindent
다음으로 $E'$를 $P$가 $E'$ 위에서 완전히 분해되게 하는
$F$의 또 다른 최소 체확대라 하자.
$P = c (x - \alpha_1) \ldots (x - \alpha_n)$으로 쓰자.
여기서 $c \in F$이고 $\alpha_i \in E'$이다. 최소성으로부터 다시
$E' = F(\alpha_1, \ldots, \alpha_n)$임이 따른다. 더욱이 임의의
$\sigma : E' \to \overline{F}$를 택하면
(Lemma \ref{fields-lemma-map-into-algebraic-closure})
어떤 순열 $\tau : \{1, \ldots, n\} \to \{1, \ldots, n\}$에 대하여
$\sigma(\alpha_i) = \beta_{\tau(i)}$임을 곧 알 수 있다.
따라서 $\sigma(E') = E$이다. Lemma
\ref{fields-lemma-characterize-normal}에 의해 이는 $E'$가 $F$의
정규확대이고, $F$의 확대들로서 $E \cong E'$임을 뜻한다.
이로써 증명이 끝난다.
\end{proof}

\begin{definition}
\label{fields-definition-splitting-field}
$F$를 체라 하고 $P \in F[x]$를 비상수다항식이라 하자.
Lemma \ref{fields-lemma-splitting-field}에서 구성한 체확대 $E/F$를
{\it $F$ 위의 $P$의 분해체}라고 한다.
\end{definition}

\begin{lemma}
\label{fields-lemma-normal-closure}
\begin{slogan}
유한 체확대의 정규폐포가 존재한다.
\end{slogan}
$E/F$를 체의 유한확대라 하자. $K$가 $F$ 위에서 정규가 되게
하는 가장 작은 유한확대 $K/E$가 유일하게 존재한다.
\end{lemma}

\begin{proof}
$\alpha_1, \ldots, \alpha_n$을 $F$ 위에서 $E$를 생성하는
원소들로 택하자. $P_1, \ldots, P_n$을
$\alpha_1, \ldots, \alpha_n$의 $F$ 위 최소다항식들이라 하자.
$P = P_1 \ldots P_n$으로 놓는다. 각 $(x - \alpha_i)$가
$P_i$를 나누므로 $(x - \alpha_1) \ldots (x - \alpha_n)$은
$P$를 나눈다. 다음과 같이 쓰자.
$P = (x - \alpha_1) \ldots (x - \alpha_n)Q$.
$K/E$를 $P$의 $E$ 위 분해체라 하자. $K$가 $P$의 $F$ 위
분해체이기도 하다고 주장한다(그러면 $K$가 $F$ 위에서
정규임이 따른다). 이는 자명하다. 실제로 $K/E$는 $E$ 위에서
$Q$의 근들로 생성되고, $E$는 $F$ 위에서
$(x - \alpha_1) \ldots (x - \alpha_n)$의 근들로 생성되므로
$K$는 $F$ 위에서 $P$의 근들로 생성된다.

\medskip\noindent
유일성. $K'/E$가 $K'/F$가 정규가 되게 하는 두 번째 최소
확대라고 하자. 대수적 폐포 $\overline{F}$와 매입
$\sigma_0 : E \to \overline{F}$를 택하자. Lemma
\ref{fields-lemma-map-into-algebraic-closure}에 의해 $\sigma_0$를
$\sigma : K \to \overline{F}$ 및
$\sigma' : K' \to \overline{F}$로 연장할 수 있다.
Lemma \ref{fields-lemma-intersect-normal}에 의해
$\sigma(K) \cap \sigma'(K')$는 $F$ 위에서 정규이다.
최소성으로부터 $\sigma(K) = \sigma'(K')$라는 결론을 얻는다.
따라서 $\sigma^{-1} \circ \sigma' : K' \to K$는 $E$의
확대들의 동형사상을 준다.
\end{proof}

\begin{definition}
\label{fields-definition-normal-closure}
$E/F$를 체의 유한확대라 하자. Lemma
\ref{fields-lemma-normal-closure}에서 구성한 체확대 $K/E$를
{\it $F$ 위의 $E$의 정규폐포}라고 한다.
% P01-E0254: frozen source line 2127 omits “of” before E.
\end{definition}

\noindent
주어진 임의의 정규확대 안에서 정규폐포를 구성할 수 있다.

\begin{lemma}
\label{fields-lemma-normal-closure-inside-normal}
$L/K$를 대수적 정규확대라 하자.
\begin{enumerate}
\item $L/M/K$가 $M/K$가 유한인 부분확대이면,
$M'/K$가 유한이고 정규인 탑 $L/M'/M/K$가 존재한다.
\item $L/M'/M/K$가 $M/K$가 정규이고 $M'/M$이 유한인 탑이면,
$M''/M$이 유한이고 $M''/K$가 정규인 탑
$L/M''/M'/M/K$가 존재한다.
\end{enumerate}
\end{lemma}

\begin{proof}
(1)의 증명. $M'$를 $M$을 포함하고 $K$ 위에서 정규인
$L/K$의 가장 작은 부분확대라 하자. Lemma
\ref{fields-lemma-normal-closure}에 의해 이는 $M/K$의 정규폐포이고
$K$ 위에서 유한이다.

\medskip\noindent
(2)의 증명. $\alpha_1, \ldots, \alpha_n \in M'$이 $M$ 위에서
$M'$을 생성한다고 하자. $P_1, \ldots, P_n$을
$\alpha_1, \ldots, \alpha_n$의 $K$ 위 최소다항식들이라 하자.
$\alpha_{i, j}$를 $P_i$의 $L$ 안의 근들이라 하자.
$M''  = M(\alpha_{i, j})$로 놓는다. Lemma
\ref{fields-lemma-normally-generated}를 생성원 집합
$M \cup \{\alpha_{i, j}\}$에 적용하면 $M''$가 $K$ 위에서
정규임이 따른다.
\end{proof}

\noindent
다음 보조정리는 때때로 정규폐포의 성질을 증명하는 데 사용할 수 있다.

\begin{lemma}
\label{fields-lemma-normal-closure-tensor-product}
$L/K$를 유한확대라 하고 $M/L$을 $K$ 위에서 $L$의
정규폐포라 하자. 그러면 $K$-대수들의 전사사상
$$
L \otimes_K L \otimes_K \ldots \otimes_K L \longrightarrow M
$$
이 존재하며, 텐서인자의 수는
$[L : K]_s \leq [L : K]$로 택할 수 있다.
\end{lemma}

\begin{proof}
$K$의 대수적 폐포 $\overline{K}$를 택하자. Lemma
\ref{fields-lemma-separable-degree}의 등식에 따라
$n = [L : K]_s = |\Mor_K(L, \overline{K})|$로 놓는다.
$\Mor_K(L, \overline{K}) = \{\sigma_1, \ldots, \sigma_n\}$이라고 하자.
$M' \subset \overline{K}$를 $\sigma_i(L)$,
$i = 1, \ldots, n$으로 생성되는 $K$-부분대수라 하자.
$\overline{K}$의 모든 $K$-부분대수는 체이므로 $M'$는 체이다.
$M'$에서 $\overline{K}$로 가는 모든 $K$-대수 사상은
$\sigma_i$들을 치환하므로 $M'$를 $M'$ 안으로, 그리고 그 위로 보낸다.
구성에 의해 체 $M'$는 $\sigma_1(L)$의 원소들의 켤레들로 생성된다.
그러므로 Lemma \ref{fields-lemma-characterize-normal}에서
$M'$가 $K$ 위에서 정규이고, $\sigma_1(L)$을 포함하는
$\overline{K}$의 가장 작은 정규 부분확대임이 따른다.
정규폐포의 유일성에 의해 $M \cong M'$이다. 끝으로 전사사상
$$
L \otimes_K L \otimes_K \ldots \otimes_K L \longrightarrow M',
\quad
\lambda_1 \otimes \ldots \otimes \lambda_n \longmapsto
\sigma_1(\lambda_1) \ldots \sigma_n(\lambda_n)
$$
이 있고, 정의에 의해 $n \leq [L : K]$임에 유의하자.
\end{proof}











\section{1의 거듭제곱근}
\label{fields-section-roots-of-1}

\noindent
$F$를 체라 하자. 정수 $n \geq 1$에 대하여
$$
\mu_n(F) = \{\zeta \in F \mid \zeta^n = 1\}
$$
로 놓는다. 이를 {\it $n$차 단위근의 군} 또는
{\it $1$의 $n$제곱근들}이라고 한다. 이는 곱셈에 관하여 가환군이고
중립원은 $1$이다. 체에서 차수가 $d$인 다항식의 근의 수는 언제나
많아야 $d$임에 유의하자. 따라서 차수가 $n$인 다항방정식으로
정의되므로 $|\mu_n(F)| \leq n$이다. 물론 $\mu_n(F)$의 모든
원소의 위수는 $n$을 나눈다. 더욱이 부분군들
$$
\mu_d(F) \subset \mu_n(F),\quad d | n
$$
각각의 원소 수는 많아야 $d$이다. 이는 $\mu_n(F)$가 순환군임을
뜻한다.

\begin{lemma}
\label{fields-lemma-cyclic}
$A$를 지수가 $n$의 약수인 가환군이라 하자. 모든 $d | n$에 대하여
$\{x \in A \mid dx = 0\}$의 기수가 많아야 $d$라고 하자.
그러면 $A$는 위수가 $n$의 약수인 순환군이다.
\end{lemma}

\begin{proof}
조건들로부터 $|A| \leq n$이므로, 특히 $A$는 유한이다.
유한가환군의 구조에 의해 어떤 정수들
$1 < e_1 | e_2 | \ldots | e_r$에 대하여
$A = \mathbf{Z}/e_1\mathbf{Z} \oplus \ldots \oplus \mathbf{Z}/e_r\mathbf{Z}$
이다. 이는 $\{x \in A \mid e_1 x = 0\}$의 기수가
$e_1^r$임을 뜻한다. 따라서 $r = 1$이다.
\end{proof}

\noindent
이를 체 $\mathbf{F}_p$에 적용하면 군
$(\mathbf{Z}/p\mathbf{Z})^*$가 순환군이라는 유명한 결과를 얻는다.
이에 관해서는 유한체 절에서 더 다룬다.

\medskip\noindent
흔히 유용한 관찰이 하나 더 있다. $F$의 표수가 $p > 0$이면
$\mu_{p^n}(F) = \{1\}$이다. Lemma
\ref{fields-lemma-nr-roots-unchanged}의 증명에서 보았듯이 표수가 $p$인
체에서 $p$제곱을 취하는 사상이 단사이기 때문이다.
(물론 그 보조정리의 명제 자체로부터도 따른다.)





\section{유한체}
\label{fields-section-finite}

\noindent
$F$를 유한체라 하자. 단사사상 $\mathbf{Q} \to F$는 있을 수
없으므로 $F$의 표수가 양수임은 자명하다. $F$의 표수를 $p$라
하자. 확대 $\mathbf{F}_p \subset F$는 유한이다. 따라서 어떤
$f \geq 1$에 대하여 $F$는 $q = p^f$개의 원소를 갖는다.

\medskip\noindent
단위원군 $F^*$를 생각하자. 이는 유한가환군이므로 어떤 지수
$e$를 갖는다. 그러면 $F^* = \mu_e(F)$이고, Section
\ref{fields-section-roots-of-1}의 논의에서 $F^*$가 위수 $q - 1$인
순환군임을 알 수 있다. (사후적으로 $e = q - 1$도 따른다.)
특히 $\alpha \in F^*$가 생성원이면
$$
F = \mathbf{F}_p(\alpha)
$$
임이 자명하다. 다시 말해 확대 $F/\mathbf{F}_p$는 한 원소로
생성된다. 물론 유한체들의 모든 확대 $E/F$도 마찬가지이다
($E$가 이미 소체 위에서 한 원소로 생성되기 때문이다).





\section{원시원소}
\label{fields-section-primitive-element}

\noindent
$E/F$를 체의 유한확대라 하자. $E = F(\alpha)$이면
원소 $\alpha \in E$를 {\it $F$ 위의 $E$의 원시원소}라고 한다.

\begin{lemma}[원시원소]
\label{fields-lemma-primitive-element}
$E/F$를 체의 유한확대라 하자. 다음은 동치이다.
\begin{enumerate}
\item $F$ 위의 $E$의 원시원소가 존재한다.
\item 부분확대 $E/K/F$가 유한개뿐이다.
\end{enumerate}
더욱이 $E/F$가 분리 가능하면 (1)과 (2)가 성립한다.
\end{lemma}

\begin{proof}
$\alpha \in E$를 원시원소라 하자. $P$를 $\alpha$의
$F$ 위 최소다항식이라 하자. 다음으로 $E/K/F$를 부분확대라 하자.
$Q$를 $\alpha$의 $K$ 위 최소다항식이라 하자.
$\deg(Q) = [E : K]$임에 유의하자.
$Q = x^d + \sum_{i < d} a_i x^i$로 쓰고,
$K$가 $L = F(a_0, \ldots, a_{d - 1})$과 같다고 주장하자.
실제로 $\alpha$의 $L$ 위 차수는 $d$이고 $L \subset K$이다.
따라서 $[E : L] = [E : K]$이고, 이로부터
$[K : L] = 1$, 즉 $K = L$이 따른다. 그러므로 다항식 $Q$의
가능성이 유한개뿐임을 보이면 충분하다. 이는 자명하다.
실제로 $K[x]$에서, 따라서 특히 $E[x]$에서 인수분해
$P = QR$가 있다. $E[x]$에서 유일인수분해가 성립하므로
$E[x]$에서 $P$의 최고차항 계수가 1인 인수는 유한개뿐이다.

\medskip\noindent
$F$가 유한체이면(동치로 $E$가 유한체이면) Section
\ref{fields-section-finite}의 논의에 의해 $E/F$는 원시원소를 갖는다.
다음으로 $F$가 무한이고 진부분체 $E/K/F$가 유한개 이하라고
가정하자. 이를 $K_1, \ldots, K_N$이라 쓰자. 각
$K_i \subset E$는 진 $F$-벡터부분공간이다. $F$가 무한이므로
모든 $i$에 대하여 $\alpha \not \in K_i$인 벡터
$\alpha \in E$를 찾을 수 있다(벡터공간은 유한개의
진벡터부분공간의 합집합과 같을 수 없다. 자세한 내용은 생략한다).
그러면 $\alpha$는 $F$ 위의 $E$의 원시원소이다.

\medskip\noindent
(1)과 (2)의 동치를 확립했으므로 이제 보조정리의 마지막 명제로
넘어가자. $F$의 대수적 폐포 $\overline{F}$를 택한다.
원소들을 $\sigma_1, \ldots, \sigma_n \in
\Mor_F(E, \overline{F})$로 열거하자. $E/F$가 분리 가능하므로
Lemma \ref{fields-lemma-separable-equality}에 의해 $n = [E : F]$이다.
$i \not = j$이면
$$
V_{ij} = \Ker(\sigma_i - \sigma_j : E \longrightarrow \overline{F})
$$
가 $E$와 같지 않음에 유의하자. 따라서 앞 단락에서와 같이
모든 $i \not = j$에 대하여 $\alpha \not \in V_{ij}$인
$\alpha \in E$를 찾을 수 있다. 이로부터
$|\Mor_F(F(\alpha), \overline{F})| \geq n$이 따른다.
한편 $[F(\alpha) : F] \leq [E : F]$이다. 따라서 Lemma
\ref{fields-lemma-separable-equality}에 의해 등식이 성립하고,
$E = F(\alpha)$라는 결론을 얻는다.
\end{proof}





\section{대각합과 노름}
\label{fields-section-trace-pairing}

\noindent
$L/K$를 체의 유한확대라 하자. Lemma
\ref{fields-lemma-vector-space-is-free}에 의해 $K$-가군의 동형사상
$L \cong K^{\oplus n}$을 택할 수 있다. 물론 $n = [L : K]$은
체확대의 차수이다. 이 동형사상을 사용하여 $K$-대수 사상
$$
L \longrightarrow \text{Mat}(n \times n, K),\quad
\alpha \longmapsto \text{matrix of multiplication by }\alpha
$$
을 얻는다. 따라서 $\alpha \in L$이 주어지면 해당 행렬의
대각합과 행렬식을 취할 수 있다. 물론 이 양들은 위에서 택한
기저의 선택과 무관하다. 더 표준적으로는 $L$을 유한차원
$K$-벡터공간으로 보기만 하면
$\text{Trace}_K(\alpha : L \to L)$와 행렬식
$\det_K(\alpha : L \to L)$을 얻는다.

\begin{definition}
\label{fields-definition-trace-norm}
$L/K$를 체의 유한확대라 하자. $\alpha \in L$에 대하여
{\it 대각합}
$\text{Trace}_{L/K}(\alpha) = \text{Trace}_K(\alpha : L \to L)$
과 {\it 노름}
$\text{Norm}_{L/K}(\alpha) = \det_K(\alpha : L \to L)$
을 정의한다.
\end{definition}

\noindent
정의에서 $\text{Trace}_{L/K}$가 $K$-선형이고
$\alpha \in K$에 대하여
$\text{Trace}_{L/K}(\alpha) = [L : K]\alpha$를 만족함은 자명하다.
마찬가지로 $\text{Norm}_{L/K}$는 곱셈적이고
$\alpha \in K$에 대하여
$\text{Norm}_{L/K}(\alpha) = \alpha^{[L : K]}$이다.
이는 Exercises의 Exercises \ref{exercises-exercise-trace-det}와
\ref{exercises-exercise-trace-det-rings}에서 논의한 더 일반적인
구성의 특수한 경우이다.
% P01-E1269: frozen source lines 2394--2396 duplicate “Exercises”.

\begin{lemma}
\label{fields-lemma-characteristic-vs-minimal-polynomial}
$L/K$를 체의 유한확대라 하자. $\alpha \in L$이고 $P$가
$\alpha$의 $K$ 위 최소다항식이라 하자. 그러면
$K$-선형사상 $\alpha : L \to L$의 특성다항식은
$P^e$와 같다. 여기서 $e \deg(P) = [L : K]$이다.
\end{lemma}

\begin{proof}
$\beta_1, \ldots, \beta_e$를 $K(\alpha)$ 위에서 $L$의 기저라
하자. 그러면 $e$는 Lemma \ref{fields-lemma-degree-minimal-polynomial}와
\ref{fields-lemma-multiplicativity-degrees}에 의해
$e \deg(P) = [L : K]$이다. 이제
$L = \bigoplus K(\alpha) \beta_i$는 $\alpha$-불변 부분공간들로
직합분해되므로, $\alpha : L \to L$의 특성다항식은
$\alpha : K(\alpha) \to K(\alpha)$의 특성다항식의
$e$제곱과 같다.

\medskip\noindent
증명을 끝내기 위해 $L = K(\alpha)$라고 가정해도 된다.
이 경우 케일리--해밀턴 정리에 의해 $\alpha$는 특성다항식의
근이다. 그리고 특성다항식과 최소다항식의 차수가 같으므로
두 다항식은 같다.
\end{proof}

\begin{lemma}
\label{fields-lemma-trace-and-norm-from-minimal-polynomial}
$L/K$를 체의 유한확대라 하자. $\alpha \in L$이고
$P = x^d + a_1 x^{d - 1} + \ldots + a_d$가 $\alpha$의
$K$ 위 최소다항식이라 하자. 그러면
$$
\text{Norm}_{L/K}(\alpha) = (-1)^{[L : K]} a_d^e
\quad\text{and}\quad
\text{Trace}_{L/K}(\alpha) = - e a_1
$$
이다. 여기서 $e d = [L : K]$이다.
\end{lemma}

\begin{proof}
Lemma \ref{fields-lemma-characteristic-vs-minimal-polynomial}와 정의들에서
즉시 따른다.
\end{proof}

\begin{lemma}
\label{fields-lemma-trace-and-norm-linear}
$L/K$를 체의 유한확대라 하자. $V$를 $L$ 위의 유한차원
벡터공간이라 하고 $\varphi : V \to V$를 $L$-선형사상이라 하자.
그러면
$$
\text{Trace}_K(\varphi : V \to V) =
\text{Trace}_{L/K}(\text{Trace}_L(\varphi : V \to V))
$$
이고
$$
\det\nolimits_K(\varphi : V \to V) =
\text{Norm}_{L/K}(\det\nolimits_L(\varphi : V \to V))
$$
이다.
\end{lemma}

\begin{proof}
$V = L^{\oplus n}$인 동형사상을 택하여 $\varphi$가
$n \times n$ 행렬에 대응하게 하자. 대각합의 경우 공식의
양변은 $\varphi$에 대하여 가법적이다. 따라서 $\varphi$가
$(i, j)$ 위치에 정확히 하나의 영이 아닌 성분만 갖는 행렬에
대응한다고 가정할 수 있다. 이 경우 직접 계산하면 양변이 같음을
알 수 있다.

\medskip\noindent
노름의 경우 $\varphi$가 영이 아닌 핵을 가지면 양변이 모두
영이다. 따라서 $\varphi$가 $\text{GL}_n(L)$의 원소에
대응한다고 가정할 수 있다. 공식의 양변은 $\varphi$에 대하여
곱셈적이다. $\text{GL}_n(L)$의 모든 원소는 기본행렬들의 곱이므로
$\varphi$가 다음 두 꼴 가운데 하나라고 가정할 수 있다.
$$
E_{12}(\lambda) =
\left(
\begin{matrix}
1 & \lambda & \ldots \\
0 & 1 & \ldots \\
\ldots & \ldots & \ldots
\end{matrix}
\right)
\quad\text{or}\quad
E_1(a) =
\left(
\begin{matrix}
a & 0 & \ldots \\
0 & 1 & \ldots \\
\ldots & \ldots & \ldots
\end{matrix}
\right)
$$
(원하면 기저원들을 치환할 수도 있기 때문이다).
두 경우 모두 직접 계산으로 공식을 쉽게 확인할 수 있다.
\end{proof}

\begin{lemma}
\label{fields-lemma-trace-and-norm-tower}
$M/L/K$를 체의 유한확대들의 탑이라 하자. 그러면
$$
\text{Trace}_{M/K} = \text{Trace}_{L/K} \circ \text{Trace}_{M/L}
\quad\text{and}\quad
\text{Norm}_{M/K} = \text{Norm}_{L/K} \circ \text{Norm}_{M/L}
$$
이다.
\end{lemma}

\begin{proof}
$M$을 $L$ 위의 벡터공간으로 보고 Lemma
\ref{fields-lemma-trace-and-norm-linear}를 적용한다.
\end{proof}

\noindent
대각합 짝지음은 대각합을 사용하여 정의한다.

\begin{definition}
\label{fields-definition-trace-pairing}
$L/K$를 체의 유한확대라 하자. $L/K$의 {\it 대각합 짝지음}은
대칭 $K$-쌍선형형식
$$
Q_{L/K} : L \times L \longrightarrow K,\quad
(\alpha, \beta) \longmapsto \text{Trace}_{L/K}(\alpha\beta)
$$
이다.
\end{definition}

\noindent
체의 유한확대가 분리 가능일 필요충분조건은 대각합 짝지음이
비퇴화인 것임이 드러난다.

\begin{lemma}
\label{fields-lemma-separable-trace-pairing}
$L/K$를 체의 유한확대라 하자. 다음은 동치이다.
\begin{enumerate}
\item $L/K$는 분리 가능이다.
\item $\text{Trace}_{L/K}$는 항등적으로 영이지 않다.
\item 대각합 짝지음 $Q_{L/K}$는 비퇴화이다.
\end{enumerate}
\end{lemma}

\begin{proof}
(3)이 (2)를 함의함은 자명하다. (2)가 성립하면
$\text{Trace}_{L/K}(\gamma) \not = 0$인 $\gamma \in L$을 택하자.
그러면 $\alpha \in L$가 영이 아닐 때
$Q_{L/K}(\alpha, \gamma/\alpha) \not = 0$이다.
따라서 $Q_{L/K}$는 비퇴화이다. 이로써 (2)와 (3)의 동치를 증명했다.

\medskip\noindent
$K$의 표수가 $p$이고, $\alpha \not \in K$,
$\alpha^p \in K$인 $L = K(\alpha)$라고 하자. 그러면
$\text{Trace}_{L/K}(1) = p = 0$이다.
$i = 1, \ldots, p - 1$에 대하여 $x^p - \alpha^{pi}$는
$\alpha^i$의 $K$ 위 최소다항식이고, Lemma
\ref{fields-lemma-trace-and-norm-from-minimal-polynomial}에 의해
$\text{Trace}_{L/K}(\alpha^i) = 0$이다.
따라서 이와 같은 순수 비분리 차수 $p$의 확대에서는
$\text{Trace}_{L/K}$가 항등적으로 영이다.

\medskip\noindent
$L/K$가 분리 가능하지 않다고 가정하자. 그러면 앞 단락에서와 같은
순수 비분리 차수 $p$의 확대 $L/K'$을 갖는 부분체
$L/K'/K$가 존재한다. Lemma \ref{fields-lemma-separable-first}와
\ref{fields-lemma-finite-purely-inseparable}를 보라. 따라서 Lemma
\ref{fields-lemma-trace-and-norm-tower}에 의해 $\text{Trace}_{L/K}$가
항등적으로 영임을 알 수 있다.

\medskip\noindent
반대로 $L/K$가 분리 가능하다고 가정하자. 차수에 관한 귀납법으로
$\text{Trace}_{L/K}$가 항등적으로 영이지 않음을 보이겠다.
따라서 Lemma \ref{fields-lemma-trace-and-norm-tower}에 의해 $L/K$가
한 원소 $\alpha$로 생성된다고 가정할 수 있다(대각합이 영이
아니면 전사임을 사용하라). 어떤 $e \geq 0$에 대하여
$\text{Trace}_{L/K}(\alpha^e)$가 영이 아님을 보여야 한다.
$P = x^d + a_1 x^{d - 1} + \ldots + a_d$를 $\alpha$의
$K$ 위 최소다항식이라 하자. 그러면 $P$는 선형사상
$\alpha : L \to L$의 특성다항식이기도 하다. Lemma
\ref{fields-lemma-characteristic-vs-minimal-polynomial}를 보라.
$L/k$가 분리 가능하므로 Lemma \ref{fields-lemma-recognize-separable}에서
$P$가 $K$의 대수적 폐포 $\overline{K}$ 안에 서로 다른
$d$개의 근 $\alpha_1, \ldots, \alpha_d$를 가짐을 알 수 있다.
% P01-E0255: frozen source line 2576 switches the base field from K to k.
따라서 이들은 $\alpha : L \to L$의 고윳값들이다.
선형대수학에 의해 $\alpha^e$의 대각합은
$\alpha_1^e + \ldots + \alpha_d^e$와 같다.
그러므로 Lemma \ref{fields-lemma-sums-of-powers}에 의해 결론을 얻는다.
\end{proof}

\noindent
$K$를 체라 하고 $Q : V \times V \to K$를 $K$ 위의 유한차원
벡터공간 위 쌍선형형식이라 하자. $\dim_K(V) = n$이라고 하자.
그러면 $Q$는 선형사상 $Q : V \to V^*$,
$v \mapsto Q(v, -)$을 정의한다. 여기서
$V^* = \Hom_K(V, K)$는 쌍대벡터공간이다. 따라서 선형사상
$$
\det(Q) : \wedge^n(V) \longrightarrow \wedge^n(V)^*
$$
을 얻는다. 기저원 $\omega \in \wedge^n(V)$를 택하면
$\det(Q)(\omega) = \lambda \omega^*$로 쓸 수 있다.
여기서 $\omega^*$는 $\wedge^n(V)^*$의 쌍대기저원이다.
$\omega$의 선택을 어떤 $c \in K^*$에 대하여 $c \omega$로
바꾸면 $\omega^*$는 $c^{-1} \omega^*$로 바뀌고, 따라서
$\lambda$는 $c^2 \lambda$로 바뀐다. 그러므로
$K/(K^*)^2$ 안의 $\lambda$의 류가 잘 정의되며, 이를
{\it $Q$의 판별식}이라 한다. 정의를 풀어 쓰면
$$
\lambda = \det(Q(v_i, v_j)_{1 \leq i, j \leq n})
$$
이다. 여기서 $\{v_1, \ldots, v_n\}$은 $K$ 위에서 $V$의
기저이다. 판별식이 영이 아닐 필요충분조건은 $Q$가 비퇴화인
것임에 유의하자.

\begin{definition}
\label{fields-definition-discriminant}
$L/K$를 체의 유한확대라 하자. {\it $L/K$의 판별식}은
대각합 짝지음 $Q_{L/K}$의 판별식이다.
\end{definition}

\noindent
위의 논의와 Lemma \ref{fields-lemma-separable-trace-pairing}에 의해
판별식이 영이 아닐 필요충분조건은 $L/K$가 분리 가능한 것임을
알 수 있다. $a \in K$에 대하여 더 정확하게는 판별식이
$K/(K^*)^2$ 안의 $a$의 류라고 말해야 하지만, 흔히
“판별식은 $a$이다”라고 말한다.

\begin{exercise}
\label{fields-exercise-quadratic-discriminant}
$L/K$를 차수가 $2$인 확대라 하자. 다음 가운데 정확히 하나만
일어남을 보여라.
\begin{enumerate}
\item 판별식이 $0$이고 $K$의 표수가 $2$이며, $L/K$는
$K$의 한 원소의 제곱근을 취하여 얻는 순수 비분리 확대이다.
\item 판별식이 $1$이고 $K$의 표수가 $2$이며, $L/K$는
차수 $2$의 분리 가능 확대이다.
\item 판별식이 제곱이 아니고 $K$의 표수가 $2$가 아니며,
$L$은 $K$에 판별식의 제곱근을 첨가하여 얻는다.
\end{enumerate}
\end{exercise}










\section{갈루아 이론}
\label{fields-section-galois-theory}

\noindent
정의는 다음과 같다.

\begin{definition}
\label{fields-definition-galois}
체확대 $E/F$가 대수적이고 분리 가능하며 정규이면
{\it 갈루아} 확대라고 한다.
\end{definition}

\noindent
유한확대가 갈루아 확대일 필요충분조건은 “올바른” 수의
자기동형사상을 갖는 것임이 드러난다.

\begin{lemma}
\label{fields-lemma-finite-Galois}
$E/F$를 체의 유한확대라 하자. $E$가 $F$ 위의 갈루아 확대일
필요충분조건은 $|\text{Aut}(E/F)| = [E : F]$인 것이다.
\end{lemma}

\begin{proof}
$|\text{Aut}(E/F)| = [E : F]$라고 가정하자. Lemma
\ref{fields-lemma-normal-and-automorphisms}에 의해 이는 $E/F$가
분리 가능하고 정규임을, 따라서 갈루아 확대임을 뜻한다.
거꾸로 $E/F$가 분리 가능하면 $[E : F] = [E : F]_s$이고,
$E/F$가 더 나아가 정규이면 Lemma
\ref{fields-lemma-normal-and-automorphisms}에 의해
$|\text{Aut}(E/F)| = [E : F]$이다.
\end{proof}

\noindent
위 보조정리에 동기를 받아 갈루아군을 다음과 같이 도입한다.

\begin{definition}
\label{fields-definition-galois-group}
$E/F$가 갈루아 확대이면 군 $\text{Aut}(E/F)$를
{\it 갈루아군}이라 하고 $\text{Gal}(E/F)$로 나타낸다.
\end{definition}

\noindent
$L/K$가 무한 갈루아 확대이면 갈루아군을 위상군으로 보아야 한다.
Section \ref{fields-section-infinite-galois}에서 이 문제로 돌아온다.

\begin{lemma}
\label{fields-lemma-galois-goes-up}
$K/E/F$를 대수적 체확대들의 탑이라 하자. $K$가 $F$ 위의
갈루아 확대이면 $K$는 $E$ 위에서도 갈루아 확대이다.
\end{lemma}

\begin{proof}
Lemma \ref{fields-lemma-normal-goes-up}와
\ref{fields-lemma-separable-goes-up}를 결합한다.
\end{proof}

\begin{lemma}
\label{fields-lemma-normal-closure-galois}
$L/K$를 체의 유한 분리 가능 확대라 하자. $M$을 $K$ 위에서
$L$의 정규폐포라 하자(Definition \ref{fields-definition-normal-closure}).
그러면 $M/K$는 갈루아 확대이다.
\end{lemma}

\begin{proof}
Lemma \ref{fields-lemma-separable-first}의 부분확대 $M/M_{sep}/K$는
Lemma \ref{fields-lemma-separable-first-normal}에 의해 정규이다.
$L/K$가 분리 가능하므로 $L \subset M_{sep}$이다. 최소성에 의해
$M = M_{sep}$이고 증명이 끝난다.
\end{proof}

\noindent
$G$를 체 $K$에 작용하는 군이라 하자(체의 자기동형사상들로
작용한다고 한다). 흔히 다음 표기를 사용할 것이다.
$$
K^G = \{x \in K \mid \sigma(x) = x \ \forall \sigma \in G\}
$$
이를 $K$ 위의 $G$의 작용에 대한 {\it 고정체}라고 한다.

\begin{lemma}
\label{fields-lemma-galois-over-fixed-field}
$K$를 체라 하고 $G$를 $K$에 충실하게 작용하는 유한군이라 하자.
그러면 확대 $K/K^G$는 갈루아 확대이고
$[K : K^G] = |G|$이며, 이 확대의 갈루아군은 $G$이다.
\end{lemma}

\begin{proof}
$\alpha \in K$가 주어졌을 때 군의 작용에 의한 $\alpha$의 궤도
$G \cdot \alpha \subset K$를 생각하자. 다항식
$$
P = \prod\nolimits_{\beta \in G \cdot \alpha} (x - \beta) \in K[x]
$$
을 생각한다. 보조정리 전체의 핵심은 이 다항식이 $G$의 작용
아래 불변이므로 $K^G$에 계수를 갖는다는 것이다.
실제로 $\tau \in G$에 대하여
$$
P^\tau = \prod\nolimits_{\beta \in G \cdot \alpha} (x - \tau(\beta)) =
\prod\nolimits_{\beta \in G \cdot \alpha} (x - \beta) = P
$$
이다. 사상 $\beta \mapsto \tau(\beta)$가 궤도
$G \cdot \alpha$의 치환이기 때문이다. 따라서 $P \in K^G[x]$이다.
또한 $\alpha$가 그 궤도의 원소이므로 $P(\alpha) = 0$이다.
그러므로 확대 $K/K^G$는 대수적이다. 더욱이 $\alpha$의
$K^G$ 위 최소다항식 $Q$는 방금 구성한 다항식 $P$를 나눈다.
따라서 $Q$는 분리 가능이고(예를 들어 Lemma
\ref{fields-lemma-recognize-separable}에 의해), $K/K^G$가 분리
가능하다는 결론을 얻는다. 그러므로 $K/K^G$는 갈루아 확대이다.
증명을 끝내려면 $[K : K^G] = |G|$임을 보이면 충분하다.
그러면 Lemma \ref{fields-lemma-finite-Galois}에 의해 $G$가 갈루아군이
되기 때문이다.

\medskip\noindent
$\alpha_i \in K$, $i = 1, \ldots, n$을 다음 성질을 갖도록
유한개 택하자. $i = 1, \ldots, n$에 대하여
$\sigma(\alpha_i) = \alpha_i$이면 $\sigma$는 $G$의 중립원이다.
다음과 같이 놓는다.
$$
L = K^G(\{\sigma(\alpha_i); 1 \leq i \leq n, \sigma \in G\}) \subset K
$$
$K$ 위의 $G$의 작용이 $L$ 위의 $G$의 작용을 유도함에 유의하자.
$L$의 $K^G$ 위 차수가 $|G|$임을 보이겠다. 이것으로 증명이
끝난다. 실제로 $L \subset K$가 진부분체이면 원소들의 목록
$\alpha_1, \ldots, \alpha_n$에 $\alpha \in K$,
$\alpha \not \in L$을 더해도 $L$이 커지지 않으므로 모순이다.
따라서 $K/K^G$가 유한인 경우로 환원되며, 이는 다음 단락에서
다룬다.

\medskip\noindent
$K/K^G$가 유한이라고 가정하자. Lemma
\ref{fields-lemma-primitive-element}에 의해 $K = K^G(\alpha)$가 되는
$\alpha \in K$를 찾을 수 있다. 이 증명의 첫 단락의 구성으로부터
$\alpha$의 $K$ 위 차수가 많아야 $|G|$임을 알 수 있다.
% P01-E0256: frozen source lines 2773--2780 gives the wrong base field K.
그러나 구성에 의해 $G$가 $K^G(\alpha) = L$에 충실하게 작용하고
Lemma \ref{fields-lemma-normal-and-automorphisms}의 부등식이 있으므로
그 차수는 $|G|$보다 작을 수 없다.
\end{proof}

\begin{theorem}[갈루아 이론의 기본정리]
\label{fields-theorem-galois-theory}
$L/K$를 갈루아군이 $G$인 유한 갈루아 확대라 하자.
그러면 $K = L^G$이고 사상
$$
\{\text{subgroups of }G\}
\longrightarrow
\{\text{subextensions }L/M/K\},\quad
H \longmapsto L^H
$$
은 전단사이다. 그 역은 $M$을 $\text{Gal}(L/M)$로 보낸다.
$G$의 정규부분군 $H$들은 정확히 $M/K$가 갈루아 확대인
부분확대 $M$들에 대응한다.
\end{theorem}

\begin{proof}
Lemma \ref{fields-lemma-galois-goes-up}에 의해 부분확대 $L/M/K$가
주어지면 확대 $L/M$은 갈루아 확대이다. 물론 $L/M$은 유한이기도
하다(Lemma \ref{fields-lemma-finite-goes-up}). 따라서 Lemma
\ref{fields-lemma-finite-Galois}에 의해
$|\text{Gal}(L/M)| = [L : M]$이다. 거꾸로 $H \subset G$가
유한부분군이면 Lemma \ref{fields-lemma-galois-over-fixed-field}에 의해
$[L : L^H] = |H|$이다. 이 두 관찰로부터 정리에 주어진 전단사
대응을 형식적으로 얻는다.

\medskip\noindent
$H \subset G$가 정규이면 $L^H$는 $G$의 작용에 의해 고정되고,
표준사상 $G/H \to \text{Aut}(L^H/K)$을 얻는다.
$\text{Gal}(L/L^H) = H$이므로 이 사상은 단사여야 한다.
따라서 $|G/H| = [L^H : K]$이고, Lemma
\ref{fields-lemma-finite-Galois}에 의해 $L^H$는 갈루아 확대이다.

\medskip\noindent
거꾸로 $K \subset M \subset L$이고 $M/K$가 갈루아 확대라고
가정하자. Lemma \ref{fields-lemma-lift-maps}에 의해 모든 원소
$\tau \in \text{Gal}(L/K)$는 원소
$\tau|_M \in \text{Gal}(M/K)$를 유도한다. 이는 갈루아군의
준동형 $\text{Gal}(L/K) \to \text{Gal}(M/K)$을 유도하며,
그 핵은 $H$이다. 따라서 $H$는 정규부분군이다.
\end{proof}

\begin{lemma}
\label{fields-lemma-ses-galois}
$L/M/K$를 체들의 탑이라 하자. $L/K$와 $M/K$가 유한 갈루아
확대라고 가정하자. 그러면 유한군들의 짧은 완전열
$$
1 \to \text{Gal}(L/M) \to \text{Gal}(L/K) \to \text{Gal}(M/K) \to 1
$$
을 얻는다.
\end{lemma}

\begin{proof}
실제로 Lemma \ref{fields-lemma-lift-maps}에 의해 모든 원소
$\tau \in \text{Gal}(L/K)$는 원소
$\tau|_M \in \text{Gal}(M/K)$를 유도하며, 이것이 오른쪽의
준동형을 준다. 왼쪽 사상은 왼쪽 군을 오른쪽 화살표의 핵과
동일시한다. 유한확대들의 탑에서 차수의 곱셈성
(Lemma \ref{fields-lemma-multiplicativity-degrees})에 의해 군들의
크기가 정확히 맞으므로 이 열은 완전하다.
오른쪽 사상이 전사임을 보이기 위해 Lemma
\ref{fields-lemma-lift-maps}를 직접 사용할 수도 있다.
\end{proof}

\section{무한 갈루아 이론}
\label{fields-section-infinite-galois}

\noindent
갈루아군에는 표준 위상이 주어진다.

\begin{lemma}
\label{fields-lemma-galois-profinite}
$E/F$를 갈루아 확대라 하자. $E$에 이산위상을 주었을 때
$$
\text{Gal}(E/F) \times E \longrightarrow E
$$
가 연속이 되도록 하는 가장 약한 위상을 $\text{Gal}(E/F)$에 주자.
그러면 다음이 성립한다.
\begin{enumerate}
\item 임의의 위상공간 $X$와 사상 $X \to \text{Gal}(E/F)$에 대하여,
작용 $X \times E \to E$가 연속이면 유도된 사상
$X \to \text{Gal}(E/F)$도 연속이다.
\item 이 위상에 의해 $\text{Gal}(E/F)$는 프로유한 위상군이 된다.
\end{enumerate}
\end{lemma}

\begin{proof}
이 증명 전체에서 $E$를 이산 위상공간으로 생각한다.
자기 사상들의 집합 $\text{Map}(E, E)$ 위의 콤팩트-열린 위상은 작용
$\text{Map}(E, E) \times E \to E$가 연속이 되게 하는 보편 위상임을
상기하자. 정확한 명제는 Topology, Example
\ref{topology-example-automorphisms-of-a-set}을 보라. 보조정리에서
$\text{Gal}(E/F)$에 준 위상은 단사사상
$\text{Gal}(E/F) \to \text{Map}(E, E)$로부터 유도되는 위상이다.
따라서 보편성 (1)은 콤팩트-열린 위상의 해당 보편성에서 따른다.
가역 자기 사상들의 집합 $\text{Aut}(E)$는 콤팩트-열린 위상을 주면
위상군을 이루고(Topology, Example
\ref{topology-example-automorphisms-of-a-set} 참조),
$\text{Gal}(E/F) = \text{Aut}(E/F) \to \text{Map}(E, E)$는
$\text{Aut}(E)$를 통해 분해되므로 위상군을 얻는다.
달리 말해 단사
$$
\text{Gal}(E/F) \subset \text{Aut}(E)
$$
를 사용하여 $\text{Gal}(E/F)$에 유도된 위상군 구조를 준 것이다
(Topology, Section \ref{topology-section-topological-groups} 참조).
구성에 의해 이는 작용 $\text{Gal}(E/F) \times E \to E$가 연속이
되게 하는 가장 약한 위상군 구조이다.

\medskip\noindent
$\text{Gal}(E/F)$가 프로유한임을 보이기 위해 다음과 같이 논증한다
(갈루아군이 유한군들의 역극한임을 보이기 전에 위상을 정의했으므로
우리의 논증은 필연적으로 표준적이지 않다).
Topology, Lemma \ref{topology-lemma-profinite-group}에 의해
$\text{Gal}(E/F)$의 바탕 위상공간이 프로유한임을 보이면 충분하다.
임의의 부분집합 $S \subset E$에 대하여 다음 집합을 생각하자.
$$
G(S) = \{ f : S \to E \mid
\begin{matrix}
f(\alpha)\text{ is a root of the minimal polynomial}\\
\text{of }\alpha\text{ over }F\text{ for all }\alpha \in S
\end{matrix}
\}
$$
다항식은 유한개의 근만 가지므로 유한인 모든 $S \subset E$에 대해
$G(S)$가 유한임을 알 수 있다. $S \subset S'$이면 제한으로 사상
$G(S') \to G(S)$를 얻는다. 또한 $\alpha \in S \cap F$이고
$f \in G(S)$이면 이 경우 최소다항식이 일차이므로
$f(\alpha) = \alpha$임을 관찰하자. 프로유한 위상공간
$$
G = \lim_{S \subset E\text{ finite}} G(S)
$$
를 생각한다. 표준사상
$$
c : \text{Gal}(E/F) \longrightarrow G,\quad
\sigma \longmapsto (\sigma|_S : S \to E)_S
$$
을 생각하자. 이 사상은 단사이고, 정의를 풀어 쓰면 위에서 정의한
$\text{Gal}(E/F)$의 위상이 $G$로부터 유도된 위상임을 알 수 있다.
원소 $(f_S) \in G$가 $c$의 상에 속할 필요충분조건은 의미가 통할
때마다(즉 $\alpha, \beta, \alpha + \beta, \alpha\beta \in S$일 때마다)
(A) $f_S(\alpha) + f_S(\beta) = f_S(\alpha + \beta)$와
(M) $f_S(\alpha)f_S(\beta) = f_S(\alpha\beta)$가 성립하는 것이다.
실제로 이는 $\lim f_S : E \to E$가 $F$-대수 사상이 됨을 뜻하고,
따라서 Lemma \ref{fields-lemma-algebraic-extension-self-map}에 의해
자기동형사상이 된다. 주어진 삼중항 $(S, \alpha, \beta)$에 대한
조건 (A)와 (M)은 $G$의 닫힌 부분집합을 정의한다. 그러므로
$\text{Gal}(E/F)$는 프로유한 공간의 닫힌 부분집합과 위상동형이고,
따라서 그 자체가 프로유한이다.
\end{proof}

\begin{lemma}
\label{fields-lemma-galois-infinite}
$L/M/K$를 체들의 탑이라 하자. $L/K$와 $M/K$가 모두 갈루아라고
가정하자. 그러면 표준적인 전사 연속 준동형
$c : \text{Gal}(L/K) \to \text{Gal}(M/K)$이 존재한다.
\end{lemma}

\begin{proof}
Lemma \ref{fields-lemma-lift-maps}에 의해 $\text{Gal}(L/K)$의 원소
$\tau : L \to L$가 주어지면 제한 $\tau|_M : M \to M$은
$\text{Gal}(M/K)$의 원소이다. 이로써 준동형 $c$가 정의된다.
연속성은 위상의 보편성에서 따른다. 즉 작용
$$
\text{Gal}(L/K) \times M \longrightarrow M,\quad
(\tau, x) \longmapsto \tau(x) = c(\tau)(x)
$$
은 $M \subset L$이고 작용 $\text{Gal}(L/K) \times L \to L$이
연속이므로 연속이다. 따라서 Lemma
\ref{fields-lemma-galois-profinite}의 (1)에 의해 $c$는 연속이다.
Lemma \ref{fields-lemma-lift-maps}는 이 사상이 전사임도 보인다.
\end{proof}

\noindent
무한 갈루아 확대의 갈루아군을 생각하는 더 표준적인 방법은 다음과 같다.

\begin{lemma}
\label{fields-lemma-infinite-galois-limit}
$L/K$를 갈루아군이 $G$인 갈루아 확대라 하자.
$\Lambda$를 유한 갈루아 부분확대들의 집합이라 하자. 즉
$\lambda \in \Lambda$는 $L/L_\lambda/K$에 대응하고,
$L_\lambda/K$는 갈루아군이 $G_\lambda$인 유한 갈루아 확대이다.
$\lambda \geq \lambda'$일 필요충분조건을
$L_\lambda \supset L_{\lambda'}$로 정하여 $\Lambda$에
부분순서를 정의하자. 그러면 다음이 성립한다.
\begin{enumerate}
\item $\Lambda$는 유향 부분순서집합이다.
\item $L_\lambda$는 $\Lambda$ 위의 $K$-확대들의 계이고
$L = \colim L_\lambda$이다.
\item $G_\lambda$는 $\Lambda$ 위의 유한군들의 역계이고,
전이사상들은 전사이며
$$
G = \lim_{\lambda \in \Lambda} G_\lambda
$$
가 프로유한군으로서 성립한다.
\item 각 사영 $G \to G_\lambda$는 연속이고 전사이다.
\end{enumerate}
\end{lemma}

\begin{proof}
$K$를 포함하는 $L$의 모든 부분체는 $K$ 위에서 분리 가능하다
(이는 정의에서 즉시 따른다). $S \subset L$을 유한 부분집합이라 하자.
그러면 $K(S)/K$는 유한이고, $E/K$가 유한 갈루아가 되는 탑
$L/E/K(S)/K$가 존재한다. Lemma
\ref{fields-lemma-normal-closure-inside-normal}를 보라.
따라서 어떤 $\lambda \in \Lambda$에 대하여 $E = L_\lambda$이다.
특히 집합 $\Lambda$는 공집합이 아니다. 또한
$\lambda_1, \lambda_2 \in \Lambda$가 주어지면 유한집합들
$S_1, S_2 \subset L$에 대하여 $L_{\lambda_i} = K(S_i)$로 쓸 수 있다
(Lemma \ref{fields-lemma-finite-finitely-generated}). 그러면
$K(S_1 \cup S_2) \subset L_\lambda$인
$\lambda \in \Lambda$가 존재한다. 따라서
$\lambda \geq \lambda_1, \lambda_2$이고 $\Lambda$는 유향이다
(Categories, Definition \ref{categories-definition-directed-system}).
마지막으로 $L$의 모든 원소는 어떤 $\lambda \in \Lambda$에 대한
$L_\lambda$에 들어가므로, Categories, Section
\ref{categories-section-directed-colimits}의 여과 여극한에 대한 기술에서
$\colim L_\lambda = L$이 따른다.

\medskip\noindent
$\Lambda$에서 $\lambda \geq \lambda'$이면 Lemma
\ref{fields-lemma-ses-galois}에 의해 표준 전사사상
$G_\lambda \to G_{\lambda'}$, $\sigma \mapsto \sigma|_{L_{\lambda'}}$를
얻는다. 따라서 전이사상들이 전사인 유한군들의 역계를 얻는다.

\medskip\noindent
$G = \text{Aut}(L/K)$임을 상기하자. Lemma
\ref{fields-lemma-galois-infinite}에 의해 $\sigma \in G$의 $L_\lambda$로의
제한 $\sigma|_{L_\lambda}$는 $G_\lambda$의 원소이다. 더욱이 이 절차는
연속 전사사상 $G \to G_\lambda$를 준다. $G_\lambda$의 역계에서
전이사상들도 제한으로 주어지므로 표준 연속사상
$$
G \longrightarrow \lim_{\lambda \in \Lambda} G_\lambda
$$
을 얻음이 분명하다. 위상군의 범주에서 극한의 정의에 의해 이는
연속이다. 이 극한들이 집합 및 위상공간의 범주로 가는 망각함자와
교환함은 Topology, Lemma
\ref{topology-lemma-topological-group-limits}를 상기하자.
한편 $L = \colim L_\lambda$이므로 역극한의 임의의 원소를 집합으로
보면 $L$의 자기동형사상을 정의함이 분명하다. 따라서 이 사상은
전단사이다. 양변의 위상이 프로유한이고 프로유한 공간 사이의
전단사 연속사상은 위상동형이므로(Topology, Lemma
\ref{topology-lemma-bijective-map}), 증명이 끝난다.
\end{proof}

\begin{theorem}[무한 갈루아 이론의 기본정리]
\label{fields-theorem-inifinite-galois-theory}
$L/K$를 갈루아 확대라 하자. $G = \text{Gal}(L/K)$를 프로유한
위상군으로 본 갈루아군이라 하자(Lemma
\ref{fields-lemma-galois-profinite}). 그러면 $K = L^G$이고 사상
$$
\{\text{closed subgroups of }G\}
\longrightarrow
\{\text{subextensions }L/M/K\},\quad
H \longmapsto L^H
$$
은 전단사이며, 그 역은 $M$을 $\text{Gal}(L/M)$으로 보낸다.
유한 부분확대 $M$은 정확히 열린 부분군 $H \subset G$에 대응한다.
$G$의 닫힌 정규부분군 $H$는 정확히 $K$ 위의 갈루아 부분확대
$M$에 대응한다.
\end{theorem}

\begin{proof}
유한 갈루아 이론의 결과(Theorem \ref{fields-theorem-galois-theory})를
더 언급하지 않고 사용하겠다. $S \subset L$을 유한 부분집합이라 하자.
$K(S) \subset E$이고 $E/K$가 유한 갈루아가 되는 탑 $L/E/K$가
존재한다. Lemma \ref{fields-lemma-normal-closure-inside-normal}를 보라.
달리 말해 $L/K$는 그 유한 갈루아 부분확대들의 합집합이다.
그러한 $E$에 대하여 Lemma \ref{fields-lemma-galois-infinite}에 의해
사상 $\text{Gal}(L/K) \to \text{Gal}(E/K)$은 전사이고 연속이다.
즉 $\text{Gal}(E/K)$의 위상이 이산이므로 그 핵은 열린집합이다.
특히 $E^{\text{Gal}(E/K)} = K$이므로 $L \setminus K$의 어떤 원소도
$\text{Gal}(L/K)$에 의해 고정되지 않음을 알 수 있다.
이로써 $L^G = K$가 증명된다.

\medskip\noindent
Lemma \ref{fields-lemma-galois-goes-up}에 의해 부분확대 $L/M/K$가 주어지면
확대 $L/M$은 갈루아이다. $G$에 주어진 위상의 정의에서 부분군
$\text{Gal}(L/M)$이 닫혀 있음은 즉시 따른다. 위의 논증을 $L/M$에
적용하면 $L^{\text{Gal}(L/M)} = M$임을 알 수 있다.

\medskip\noindent
거꾸로 $H \subset G$를 닫힌 부분군이라 하자. 다음을 보이겠다.
$H = \text{Gal}(L/L^H)$. 포함관계
$H \subset \text{Gal}(L/L^H)$는 분명하다.
$g \in \text{Gal}(L/L^H)$라고 하자. $S \subset L$을 유한
부분집합이라 하자. $g$의 열린근방
$U_S(g) = \{g' \in G \mid g'(s) = g(s)\}$가 $H$와 만남을 보이겠다.
그러면 $H$가 닫혀 있으므로 $g \in H$이다. 증명의 첫 단락에서처럼
$K(S)$를 포함하는 유한 갈루아 부분확대 $L/E/K$를 택하고 준동형
$c : \text{Gal}(L/K) \to \text{Gal}(E/K)$를 생각하자.
그러면 $L^H \cap E = E^{c(H)}$이다. $g$는 $L^H$를 고정하므로
$E^{c(H)}$를 고정하고, 따라서 유한 갈루아 이론에 의해
$c(g) \in c(H)$이다. $c(h) = c(g)$인 $h \in H$를 택하자.
그러면 원하는 대로 $h \in U_S(g)$이다.

\medskip\noindent
이로써 닫힌 부분군들과 부분확대들 사이의 대응을 확립했다.

\medskip\noindent
$H \subset G$가 열린집합이라고 가정하자. 위와 같이 논증하면
충분히 큰 유한 갈루아 부분확대 $E$에 대해 $H$가
$\text{Gal}(L/E)$를 포함함을 알 수 있다. 따라서 $L^H$는 $E$에
포함되고, 이에 따라 $K$ 위에서 유한이다. 거꾸로 $M$이 유한
부분확대이면 $M$은 유한 부분집합 $S$에 의해 생성되며, 대응하는 부분군은
중립원 $e \in G$에 대한 열린집합 $U_S(e)$이다.

\medskip\noindent
$K \subset M \subset L$이고 $M/K$가 갈루아라고 가정하자.
Lemma \ref{fields-lemma-galois-infinite}에 의해 갈루아군들의 전사 연속 준동형
$\text{Gal}(L/K) \to \text{Gal}(M/K)$이 존재하며 그 핵은
$\text{Gal}(L/M)$이다. 따라서 $\text{Gal}(L/M)$은 닫힌 정규부분군이다.

\medskip\noindent
마지막으로 $N \subset G$가 닫힌 정규부분군이라고 가정하자.
증명의 첫 단락에서와 같은 임의의 $L/E/K$에 대하여 상
$c(N) \subset \text{Gal}(E/K)$은 정규부분군이다. 그러므로
$L^N = \bigcup E^{c(N)}$은 $K$의 갈루아 확대들의 합집합이고
(유한 갈루아 이론에 의해), 따라서 $K$ 위에서 갈루아이다.
\end{proof}

\begin{lemma}
\label{fields-lemma-ses-infinite-galois}
$L/M/K$를 체들의 탑이라 하자. $L/K$와 $M/K$가 갈루아라고 가정하자.
그러면 프로유한 위상군들의 짧은 완전열
$$
1 \to \text{Gal}(L/M) \to \text{Gal}(L/K) \to \text{Gal}(M/K) \to 1
$$
을 얻는다.
\end{lemma}

\begin{proof}
이는 Lemma \ref{fields-lemma-galois-infinite}를 다시 표현한 것이다.
\end{proof}

\section{복소수}
\label{fields-section-complex-numbers}

\noindent
대수학의 기본정리는 복소수체가 대수적으로 닫힌 체라고 말한다.
이 절에서는 이를 간단히 논의한다.

\medskip\noindent
먼저 지적할 것은 이를 증명하려면 미적분학의 결과를 조금 사용해야
한다는 점이다. 실수 위의 모든 홀수차 다항식은 실근을 갖는다는
직관적으로 분명한 사실을 사용하겠다. 즉 어떤 $k \geq 0$과
$a_{2k + 1} \not = 0$에 대하여
$P(x) = a_{2k + 1} x^{2k + 1} + \ldots + a_0 \in \mathbf{R}[x]$라 하자.
$a_{2k + 1} > 0$이라고 가정해도 좋고 그렇게 하겠다.
$x \in \mathbf{R}$가 매우 크고 양수이면 항
$a_{2k + 1} x^{2k + 1}$이 다른 모든 항을 지배하므로
$P(x) > 0$이다. 마찬가지로 $x \ll 0$이면 같은 이유로
$P(x) < 0$이다(여기서 차수가 홀수라는 사실을 사용한다).
따라서 중간값정리에 의해 $P(x) = 0$인 $x \in \mathbf{R}$가 존재한다.

\medskip\noindent
위에서 얻을 수 있는 한 결론은 $\mathbf{R}$에는 자명하지 않은 홀수차
체확대가 없다는 것이다. 그러한 확대의 원소들은 홀수차 최소다항식을
가질 것이기 때문이다.

\medskip\noindent
이제 $K/\mathbf{R}$를 갈루아군이 $G$인 유한 갈루아 확대라 하자.
$P \subset G$를 $2$-Sylow 부분군이라 하자. 그러면
$K^P/\mathbf{R}$는 홀수차 확대이므로 위에 의해
$K^P = \mathbf{R}$이고, 이는 다시 $G = P$를 함의한다.
(물론 이 모든 논증은 갈루아 이론에 의존한다.) 따라서 $G$는
$2$-군이다. $G$가 자명하지 않으면 $\mathbf{C}$가 동형을 제외하면
$\mathbf{R}$의 유일한 차수 $2$ 확대이므로 $\mathbf{C} \subset K$이다.
$G$가 $2$개보다 많은 원소를 가지면 $\mathbf{C}$의 이차확대를 얻게
된다. 모든 복소수가 제곱근을 가지므로 이는 모순이다.

\medskip\noindent
결론은 $\mathbf{C}$가 대수적으로 닫혀 있다는 것이다. 실제로 그렇지
않다면 자명하지 않은 유한확대 $K/\mathbf{C}$를 얻게 되고, Lemma
\ref{fields-lemma-normal-closure}에 의해 이를 $\mathbf{R}$ 위에서 정규인
(따라서 갈루아인) 확대라고 가정할 수 있다. 그러나 위에서 이 경우
$K = \mathbf{C}$임을 보았다.

\begin{lemma}[대수학의 기본정리]
\label{fields-lemma-C-algebraically-closed}
체 $\mathbf{C}$는 대수적으로 닫혀 있다.
\end{lemma}

\begin{proof}
위의 논의를 보라.
\end{proof}

\section{쿠머 확대}
\label{fields-section-Kummer}

\noindent
$K$를 체라 하자. $n \geq 2$를 $K$가 $1$의 원시 $n$제곱근을
포함하는 정수라 하자. % P01-E1271
$a \in K$라 하자. $L$을 방정식 $x^n = a$의 근 $b$를 첨가하여 얻는
$K$의 확대라 하자. 그러면 $L/K$는 갈루아 확대이다.
$G = \text{Gal}(L/K)$가 갈루아군이면 사상
$$
G \longrightarrow \mu_n(K),\quad \sigma \longmapsto \sigma(b)/b
$$
은 군들의 단사 준동형이다. 특히 $G$는 순환군 $\mu_n(K)$의 부분군으로서
그 위수가 $n$을 나누는 순환군이다. 쿠머 이론은 그 역을 준다.

\begin{lemma}[쿠머 확대]
\label{fields-lemma-Kummer}
$L/K$를 갈루아군이 $\mathbf{Z}/n\mathbf{Z}$인 체의 갈루아 확대라 하자.
더 나아가 $K$의 표수가 $n$과 서로소이고 $K$가 $1$의 원시
$n$제곱근을 포함한다고 가정하자. 그러면 $L = K[z]$이고
$z^n \in K$이다.
\end{lemma}

\begin{proof}
$\zeta \in K$를 $1$의 원시 $n$제곱근이라 하자.
$\sigma$를 $\text{Gal}(L/K)$의 생성원이라 하자.
$\sigma : L \to L$을 $K$-선형 연산자로 생각한다.
선형 연산자로서 $\sigma^n - 1 = 0$임에 유의하자.
지표들의 선형독립성(Lemma \ref{fields-lemma-independence-characters})을
적용하면 $\sigma$를 소멸시키는 차수 $< n$인 $K$ 위의 다항식이
존재할 수 없음을 알 수 있다. 따라서 선형 연산자로서 $\sigma$의
최소다항식은 $x^n - 1$이다. $\zeta$는 $x^n - 1$의 근이므로
선형대수학에 의해 $\sigma(z) = \zeta z$인 $0 \neq z \in L$이
존재한다. 이 $z$는
$\sigma(z^n) = (\zeta z)^n = z^n$이므로 $z^n \in K$를 만족한다.
더 나아가
$z, \sigma(z), \ldots, \sigma^{n - 1}(z) =
z, \zeta z, \ldots \zeta^{n - 1} z$는 서로 다르므로 $z$가 $K$ 위에서
$L$을 생성함이 보장된다. 따라서 원하는 대로 $L = K[z]$이다.
\end{proof}

\begin{lemma}
\label{fields-lemma-adjoint-pth-root-unity}
$K$를 대수적 폐포가 $\overline{K}$인 체라 하자. $p$를 $K$의 표수와
다른 소수라 하자. $\zeta \in \overline{K}$를 $1$의 원시
$p$제곱근이라 하자. 그러면 $K(\zeta)/K$는 차수가 $p - 1$을 나누는
갈루아 확대이다.
\end{lemma}

\begin{proof}
다항식 $x^p - 1$의 근들은
$1, \zeta, \zeta^2, \ldots, \zeta^{p - 1}$이므로 이 다항식은
$K(\zeta)$ 위에서 완전히 분해된다. 따라서 $K(\zeta)/K$는 분해체이고
이에 따라 정규이다. $x^p - 1$이 분리 가능 다항식이므로 이 확대는
분리 가능하다. 따라서 이 확대는 갈루아이다. $K$ 위의
$K(\zeta)$의 임의의 자기동형사상은 어떤 $1 \leq i \leq p - 1$에
대하여 $\zeta$를 $\zeta^i$로 보낸다. 따라서 갈루아군은
$(\mathbf{Z}/p\mathbf{Z})^*$의 부분군이다.
\end{proof}

\begin{lemma}
\label{fields-lemma-subfields-kummer}
\begin{reference}
\cite[Theorem 5.2]{Radical}
\end{reference}
$K$를 체라 하자. $L/K$를 $a = \alpha^e \in K$인 원소
$\alpha$에 의해 생성되는 차수 $e$의 유한확대라 하자. $L$에 있는
1의 모든 $e$제곱근이 $K$에 포함되면,
% P01-E0258
임의의 부분확대 $L/L'/K$는 어떤 $d | e$에 대하여
$\alpha^d$에 의해 생성된다.
\end{lemma}

\begin{proof}
$d | e$에 대하여 부분체 $K(\alpha^d)$는
$[K(\alpha^d) : K] = e/d$와 $[L : K(\alpha^d)] = d$를 만족함을
관찰하자. $L/L'/K$를 부분확대라 하고 $d = [L : L']$이라 하자.
$\alpha^d \in L'$이면 차수 때문에 $L' = K(\alpha^d)$이다.
$P \in L'[x]$를 $L'$ 위의 $\alpha$의 최소다항식이라 하자.
그러면 $P$는 $x^e - a$를 나누고 $P$의 차수가 $d$이다. $L$ 위에서
$x^e - a$의 분해체 안에서 다음과 같이 쓰자.
$$
x^e - a = \prod\nolimits_{i = 1, \ldots, e} (x - \zeta_i \alpha)
$$
$\zeta_i$들은 1의 $e$제곱근들이고, 번호를 다시 매기면
$$
P = \prod\nolimits_{i = 1, \ldots, d} (x - \zeta_i \alpha)
$$
이다. $P$의 상수항은
% P01-E0259
$$
c = (\prod\nolimits_{i = 1, \ldots, d} \zeta_i) \alpha^d
$$
와 같고 $L' \subset L$에 속한다. $\alpha \in L$이므로
$\zeta = \prod_{i = 1, \ldots, d} \zeta_i$는 $L$에 속하고,
가정에 의해 $K$에 속한다. 따라서
$\alpha^d = \zeta^{-1}c \in L'$이고 결론을 얻는다.
\end{proof}

\section{아르틴--슈라이어 확대}
\label{fields-section-Artin-Schreier}

\noindent
$K$를 표수가 $p > 0$인 체라 하자. $a \in K$라 하자. $L$을 방정식
$x^p - x = a$의 근 $b$를 첨가하여 얻는 $K$의 확대라 하자.
그러면 $L/K$는 갈루아 확대이다. $G = \text{Gal}(L/K)$가 갈루아군이면
사상
$$
G \longrightarrow \mathbf{Z}/p\mathbf{Z},\quad
\sigma \longmapsto \sigma(b) - b
$$
은 군들의 단사 준동형이다. 특히 $G$는
$\mathbf{Z}/p\mathbf{Z}$의 부분군으로서 위수가 $p$를 나누는
순환군이다. 아르틴--슈라이어 확대 이론은 그 역을 준다.

\begin{lemma}[아르틴--슈라이어 확대]
\label{fields-lemma-Artin-Schreier}
$L/K$를 표수가 $p > 0$이고 갈루아군이
$\mathbf{Z}/p\mathbf{Z}$인 체들의 갈루아 확대라 하자.
그러면 $L = K[z]$이고 $z^p - z \in K$이다.
\end{lemma}

\begin{proof}
$\sigma$를 $\text{Gal}(L/K)$의 생성원이라 하자.
$\sigma : L \to L$을 $K$-선형 연산자로 생각한다.
선형 연산자로서 $\sigma^p - 1 = 0$임을 관찰하자.
지표들의 선형독립성(Lemma \ref{fields-lemma-independence-characters})을
적용하면 $\sigma$를 소멸시키는 차수 $< p$인 다항식이 존재할 수 없다.
따라서 $\sigma$의 최소다항식은 $x^p - 1 = (x - 1)^p$이다.
이는 $(\sigma - 1)^{p - 1}(w) = y$가 영이 아닌 원소가 되는
$w \in L$이 존재함을 함의한다. 그러면 $\sigma(y) = y$, 즉
$y \in K$이다. 따라서
$z = y^{-1}(\sigma - 1)^{p - 2}(w)$는
$\sigma(z) = z + 1$을 만족한다. $z \not \in K$이므로
$L = K[z]$이다. 더 나아가
$\sigma(z^p - z) = (z + 1)^p - (z + 1) = z^p - z$이므로
$z^p - z \in K$이고 증명이 끝난다.
\end{proof}

\section{초월성}
\label{fields-section-transcendence}

\noindent
표준 정의들을 상기한다.

\begin{definition}
\label{fields-definition-transcendence}
$K/k$를 체확대라 하자.
\begin{enumerate}
\item $K$의 원소들의 모임 $\{x_i\}_{i \in I}$를 생각하자.
$X_i$를 $x_i$로 보내는 사상
$$
k[X_i; i\in I] \longrightarrow K
$$
이 단사이면 이 모임은 $k$ 위에서 {\it 대수적으로 독립}이라고 한다.
\item 다항식환 $k[x_i; i \in I]$의 분수체를
$k(x_i; i\in I)$로 나타낸다.
\item $k$의 {\it 순수 초월 확대}란 $k$ 위의 다항식환의 분수체와
동형인 임의의 체확대 $K/k$이다.
\item $K/k$의 {\it 초월기저}란 $k$ 위에서 대수적으로 독립이고
확대 $K/k(x_i; i\in I)$가 대수적인 원소들의 모임
$\{x_i\}_{i \in I}$이다.
\end{enumerate}
\end{definition}

\begin{example}
\label{fields-example-pi-transcendental}
$\pi$는 유리수 계수를 갖는 영이 아닌 다항식의 근이 아니므로
체 $\mathbf{Q}(\pi)$는 순수 초월 확대이다. 특히
$\mathbf{Q}(\pi) \cong \mathbf{Q}(x)$이다.
\end{example}

\begin{lemma}
\label{fields-lemma-transcendence-degree}
$E/F$를 체확대라 하자. $F$ 위의 $E$의 초월기저가 존재한다.
임의의 두 초월기저는 같은 기수를 갖는다.
\end{lemma}

\begin{proof}
$A$를 $E$의 대수적으로 독립인 부분집합이라 하자. $G$를 $A$를
포함하고 $E/F$를 생성하는 $E$의 부분집합이라 하자.
$A \subset B \subset G$인 초월기저 $B$를 찾을 수 있음을 보이겠다.
이를 위해 $A$를 포함하는 $G$의 부분집합들 가운데 대수적으로
독립인 것들의 모임 $\mathcal{B}$를 생각한다. 포함관계로
$\mathcal{B}$에 부분순서를 정의한다. 그러면 $\mathcal{B}$는 적어도
한 원소 $A$를 포함한다. $\mathcal{B}$의 전순서 부분집합 $T$의
원소들의 합집합은 $F$ 위의 $E$의 대수적으로 독립인 부분집합이다.
실제로 대수적 종속관계가 있다면 그것은 $T$의 어떤 원소에서 이미
나타났을 것이다(다항식은 유한개의 변수만 사용하기 때문이다).
이 합집합은 또한 $A$를 포함하고 $G$에 포함된다. 조른의 보조정리에
의해 극대원소 $B \in \mathcal{B}$가 존재한다. 이제 $E$가
$F(B)$ 위에서 대수적임을 보이겠다. 그렇지 않다면
$F(G) = E$이므로 $F(B)$ 위에서 초월적인 원소 $f \in G$가 존재할
것이다. 그러면 $B \cup\{f\}$가 대수적으로 독립이 되어 $B$의
극대성에 모순이다. 따라서 $B$가 원하는 초월기저이다.

\medskip\noindent
$B$와 $B'$를 두 초월기저라 하자. 일반성을 잃지 않고
$|B'| \leq |B|$라고 가정할 수 있다. 이제 증명을 두 경우로 나눈다.
첫째로 $B$가 무한집합이라고 하자. 각 $\alpha \in B'$에 대하여
대수적 종속관계에는 유한개의 부정원만 쓰이므로 $\alpha$가
$F(B_{\alpha})$ 위에서 대수적이 되는 유한집합
$B_{\alpha} \subset B$가 존재한다. 다음과 같이 정의한다.
$B^* = \bigcup_{\alpha\in B'} B_{\alpha}$.
구성에 의해 $B^* \subset B$이지만, 실제로 두 집합이 같음을
보이겠다. 같지 않다고 하고 원소
$\beta \in B \setminus B^*$가 있다고 하자. $\beta$는
$F(B')$ 위에서 대수적이고, 후자는 $F(B^*)$ 위에서 대수적이다.
따라서 $\beta$는 $F(B^*)$ 위에서 대수적이며 이는 모순이다.
그러므로 $|B| \leq |\bigcup_{\alpha \in B'} B_{\alpha}|$이다.
이제 $B'$가 유한이면 $B$도 유한이므로 $B'$가 무한이라고 가정할
수 있다. 따라서
$$
|B| \leq |\bigcup\nolimits_{\alpha \in B'} B_{\alpha}| = |B'|
$$
이다. 각 $B_\alpha$가 유한이고 $B'$가 무한이기 때문이다.
그러므로 무한인 경우 $|B| = |B'|$이다.

\medskip\noindent
이제 $B$가 유한인 경우를 보아야 한다. 이 경우 $B'$도 유한이므로
$B = \{\alpha_1, \ldots, \alpha_n\}$이고
$B' = \{\beta_1, \ldots, \beta_m\}$이며 $m \leq n$이라고 하자.
$m$에 대한 귀납법을 사용한다. $m = 0$이면 $E/F$가 대수적이므로
$B = \emptyset$이고 따라서 $n = 0$이다. $m > 0$이면
$f$에 $x$가 실제로 나타나고
$f(\beta_1, \alpha_1, \ldots, \alpha_n) = 0$인 기약다항식
$f \in F[x, y_1, \ldots, y_n]$이 존재한다. $\beta_1$은 $F$ 위에서
대수적이지 않으므로 $f$에는 어떤 $y_i$가 나타나야 한다.
일반성을 잃지 않고 $f$에 $y_1$이 나타난다고 가정한다.
$B^* = \{\beta_1, \alpha_2, \ldots, \alpha_n\}$이라 하자.
$B^*$가 $E/F$의 기저임을 보이겠다. 이를 위해
$\alpha_1$이 $F(B^*)$ 위에서 대수적이므로 대수적 확대들의 탑
$$
E/ F(B^*, \alpha_1) / F(B^*)
$$
을 얻음을 관찰한다. 이제 $B^*$가(원소들의 중복도를 세어)
$F$ 위에서 대수적으로 독립임을 보이겠다. 그렇지 않다면
$g(\beta_1, \alpha_2, \ldots, \alpha_n) = 0$이고
$g\in F[x, y_2, \ldots, y_n]$인 기약다항식이 존재할 것이다.
이 다항식에는 $x$가 나타나야 하므로 $\beta_1$은
$F(\alpha_2, \ldots, \alpha_n)$ 위에서 대수적이다. 그러면
$\alpha_1$도 $F(\alpha_2, \ldots, \alpha_n)$ 위에서 대수적이 되며,
이는 불가능하다. 따라서 $\{\alpha_2, \ldots, \alpha_n\}$와
$\{\beta_2, \ldots, \beta_m\}$는 $F(\beta_1)$ 위의 $E$의
기저들이다. 귀납가정에 의해 $m = n$이다.
\end{proof}

\begin{definition}
\label{fields-definition-transcendence-degree}
$K/k$를 체확대라 하자. $k$ 위의 $K$의 {\it 초월차수}란
$k$ 위의 $K$의 초월기저의 기수이다. 이를 $\text{trdeg}_k(K)$로
나타낸다.
\end{definition}

\begin{lemma}
\label{fields-lemma-transcendence-degree-tower}
$L/K/k$를 체확대들이라 하자. 그러면
$$
\text{trdeg}_k(L) =
\text{trdeg}_K(L) +
\text{trdeg}_k(K).
$$
\end{lemma}

\begin{proof}
$A \subset K$를 $k$ 위의 $K$의 초월기저로 택한다.
$B \subset L$을 $K$ 위의 $L$의 초월기저로 택한다.
그러면 $A \cup B$가 $k$ 위의 $L$의 초월기저임을 쉽게 알 수 있다.
\end{proof}

\begin{example}
\label{fields-example-pi-e-transcendental}
수 $e$와 $\pi$를 첨가하여 만든 체확대 $\mathbf{Q}(e, \pi)$를
생각하자. $e$와 $\pi$가 모두 유리수체 위에서 초월적이므로 이
체확대의 초월차수는 적어도 $1$이다. 그러나 $e$와 $\pi$가
대수적으로 독립이면 이 체확대의 초월차수는 $2$일 수 있다.
이 명제가 참인지 여부는 알려져 있지 않으며, 따라서
$\text{trdeg}(\mathbf{Q}(e, \pi))$를 결정하는 문제는 미해결이다.
\end{example}

\begin{example}
\label{fields-example-function-field-not-unique-transcendence-basis}
$F$를 체라 하고 $E = F(t)$라 하자. $E = F(t)$이므로
$\{t\}$는 초월기저이다. 그러나 $F(t)/F(t^2)$가 대수적이므로
$\{t^2\}$도 초월기저이다. 이는 확대 $E/F$를 대수적 확대
$E/F'$와 순수 초월 확대 $F'/F$로 언제나 분해할 수 있지만,
그 분해는 유일하지 않고 초월기저의 선택에 의존함을 보여 준다.
\end{example}

\begin{example}
\label{fields-example-riemann-surface-transcendence}
$X$를 콤팩트 리만곡면이라 하자. 그러면 함수체 $\mathbf{C}(X)$
(Example \ref{fields-example-field-of-meromorphic-functions} 참조)는
$\mathbf{C}$ 위에서 초월차수 1을 갖는다. 실제로 $\mathbf{C}$의
초월차수 1인 {\it 임의의} 유한생성 확대는 리만곡면에서 나온다.
더 나아가 콤팩트 리만곡면들과 상수가 아닌 정칙사상들의 범주와,
$\mathbf{C}$의 초월차수 $1$인 유한생성 확대들과
$\mathbf{C}$-대수 사상들의 범주의 반대범주는 서로 동치이다.
\cite{Forster}를 보라.

\medskip\noindent
위 명제에는 대수적인 판본도 있다. 대수적으로 닫힌 체 $k$(예를 들어
복소수체) 위의 사영공간 안에 있는 (기약) 대수곡선이 주어지면 그
``유리함수체''를 생각할 수 있다. 대략 말해 이는 분모가 곡선 위에서
항등적으로 영이 아닌 다항식들의 몫처럼 보이는 함수들이다.
매끄러운 사영곡선들과 상수가 아닌 곡선 사상들의 범주와,
$k$의 초월차수 1인 유한생성 확대들의 범주 사이에도 비슷한 반동치가
있다(Algebraic Curves, Theorem
\ref{curves-theorem-curves-rational-maps}). \cite{H}를 보라.
\end{example}

\begin{definition}
\label{fields-definition-algebraically-closed-in}
$K/k$를 체확대라 하자.
\begin{enumerate}
\item {\it $K$에서 $k$의 대수적 폐포}란 $k$ 위에서 대수적인
$K$의 원소들로 이루어진 $K$의 부분체 $k'$이다.
\item $k$ 위에서 대수적인 $K$의 모든 원소가 $k$에 포함되면
$k$가 {\it $K$에서 대수적으로 닫혀 있다}고 한다.
\end{enumerate}
\end{definition}

\begin{lemma}
\label{fields-lemma-purely-transcendental-degree}
$k'/k$를 체의 유한확대라 하자. 순수 초월 확대들 사이에 유도된 확대를
생각하자. 이 확대는
$k'(x_1, \ldots, x_r)/k(x_1, \ldots, x_r)$이다. 그러면
$[k'(x_1, \ldots, x_r) : k(x_1, \ldots, x_r)] = [k' : k] < \infty$이다.
\end{lemma}

\begin{proof}
확대 차수의 곱셈성(Lemma \ref{fields-lemma-multiplicativity-degrees})에 의해
$k'$가 $k$ 위에서 한 원소 $\alpha \in k'$에 의해 생성되는 경우를
증명하면 충분하다. $f \in k[T]$를 $k$ 위의 $\alpha$의 최소다항식이라
하자. 그러면 $k'(x_1, \ldots, x_r)$는 $k$ 위에서
$\alpha, x_1, \ldots, x_r$에 의해 생성되고, 따라서
$k'(x_1, \ldots, x_r)$는 $k(x_1, \ldots, x_r)$ 위에서
$\alpha$에 의해 생성된다.
그러므로 $f$가 여전히 $k(x_1, \ldots, x_r)[T]$의 원소로서
기약임을 보이면 충분하다. 증명을 개략적으로만 설명한다.
$f$는 $T$에 관한 일계수 다항식이고 $k[x_1, \ldots, x_r, T]$에서
가능한 인수분해가 있으면 $x_i$를 $0$으로 놓아 $k[T]$에서의
인수분해를 얻게 되므로, $f$가
$k[x_1, \ldots, x_r, T]$의 원소로서 기약임은 분명하다.
가우스의 보조정리로 결론을 얻는다.
\end{proof}

\begin{lemma}
\label{fields-lemma-algebraic-closure-in-finitely-generated}
$K/k$를 유한생성 체확대라 하자. $K$에서 $k$의 대수적 폐포는
$k$ 위에서 유한이다.
\end{lemma}

\begin{proof}
$x_1, \ldots, x_r \in K$를 $k$ 위의 $K$의 초월기저라 하자.
그러면 $n = [K : k(x_1, \ldots, x_r)] < \infty$이다.
$k'/k$가 유한인 $k \subset k' \subset K$를 가정하자. 이 경우
$[k'(x_1, \ldots, x_r) : k(x_1, \ldots, x_r)] = [k' : k] < \infty$이다.
Lemma \ref{fields-lemma-purely-transcendental-degree}를 보라. 따라서
$$
[k' : k] = [k'(x_1, \ldots, x_r) : k(x_1, \ldots, x_r)] \leq
[K : k(x_1, \ldots, x_r)] = n.
$$
달리 말해 유한 부분확대들의 차수는 유계이고, 보조정리가 따른다.
\end{proof}

\section{선형적으로 서로소인 확대}
\label{fields-section-linearly-disjoint}

\noindent
$k$를 체라 하고 $K$와 $L$을 $k$의 체확대들이라 하자.
또한 $K$와 $L$이 어떤 더 큰 체 $\Omega$에 매입되어 있다고 가정하자.

\begin{definition}
\label{fields-definition-compositum}
체확대들의 도표
\begin{equation}
\label{fields-equation-inside-omega}
\vcenter{
\xymatrix{
L \ar[r] & \Omega \\
k \ar[r] \ar[u] & K \ar[u]
}
}
\end{equation}
을 생각하자. $K$와 $L$을 모두 포함하는 가장 작은 $\Omega$의 부분체를
{\it $\Omega$ 안에서 $K$와 $L$의 합성체}라 하고,
$KL$로 나타낸다.
\end{definition}

\noindent
$KL$은 $k$ 위에서 집합 $K \cup L$에 의해 생성되고, $L$ 위에서
집합 $K$에 의해 생성되며, $K$ 위에서 집합 $L$에 의해 생성됨이 분명하다.

\medskip\noindent
% P01-E0260
{\bf 주의:} 합성체의 (동형류는) $K$와 $L$을 $\Omega$에 매입하는
선택에 의존한다. 예를 들어 수체
$K = \mathbf{Q}(2^{1/8}) \subset \mathbf{R}$와
$L = \mathbf{Q}(2^{1/12}) \subset \mathbf{R}$를 생각하자.
$\mathbf{R}$ 안에서의 합성체는 $\mathbf{Q}$ 위에서 차수가 $24$인 체
$\mathbf{Q}(2^{1/24})$이다. 그러나 $x$를
$2^{1/8}e^{2\pi i/8}$로 보내어
$K = \mathbf{Q}[x]/(x^8 - 2)$를 $\mathbf{C}$에 매입하면 합성체
$\mathbf{Q}(2^{1/12}, 2^{1/8}e^{2\pi i/8})$는
$i = e^{2\pi i/4}$를 포함하고 $\mathbf{Q}$ 위에서 차수 $48$을 갖는다
(차수가 $48$임을 보이는 것은 생략하지만, $i$의 존재만으로도 두 합성체가
동형이 아님이 확실히 증명된다).

\begin{definition}
\label{fields-definition-linearly-disjoint}
(\ref{fields-equation-inside-omega})와 같은 체들의 도표를 생각하자. 사상
$$
K \otimes_k L \longrightarrow KL,\quad
\sum x_i \otimes y_i \longmapsto \sum x_i y_i
$$
이 단사이면 $K$와 $L$이 {\it $\Omega$ 안에서 $k$ 위에 선형적으로
서로소}라고 한다.
\end{definition}

\noindent
다음 보조정리는 다른 적당한 위치에 들어맞지 않는 듯하다.

\begin{lemma}
\label{fields-lemma-normal-case}
$E/F$를 정규 대수적 체확대라 하자. 다음을 만족하는 부분확대
$E / E_{sep} /F$와 $E / E_{insep} / F$가 존재한다.
\begin{enumerate}
\item $F \subset E_{sep}$는 갈루아이고 $E_{sep} \subset E$는
순수 비분리 확대이다.
\item $F \subset E_{insep}$는 순수 비분리 확대이고
$E_{insep} \subset E$는 갈루아이다.
\item $E = E_{sep} \otimes_F E_{insep}$이다.
\end{enumerate}
\end{lemma}

\begin{proof}
부분체 $E_{sep}$는 Lemma \ref{fields-lemma-separable-first}에서 얻었다.
$E_{insep} = E^{\text{Aut}(E/F)}$로 놓는다. 세부사항은 생략한다.
\end{proof}

\section{복습}
\label{fields-section-algebraic}

\noindent
이 절에서는 위에서 전개된 내용을 간단히 복습한다.

\medskip\noindent
$K/k$를 체확대라 하고 $\alpha \in K$라 하자. 그러면 다음 두 경우가
가능하다.
\begin{enumerate}
\item 원소 $\alpha$는 $k$ 위에서 초월적이다.
\item 원소 $\alpha$는 $k$ 위에서 대수적이다.
% P01-E0261
그 {\it 최소다항식}을 $P(T) \in k[T]$로 나타내자. 이는 계수가
$k$에 속하는 모닉 다항식
$P(T) = T^d + a_1 T^{d - 1} + \ldots + a_d$이다. 이 다항식은
기약이고 $P(\alpha) = 0$을 만족한다. 이러한 성질들은 $P$를 유일하게
결정하며, 정수 $d$를 {\it $\alpha$의 $k$ 위 차수}라 한다. 다시 두
경우로 나뉜다.
\begin{enumerate}
\item 다항식 $\text{d}P/\text{d}T$가 항등적으로 영이 아니다. 이는
$k$의 대수적 폐포에 있는 서로 다른 원소
$\alpha_1, \ldots, \alpha_d$에 대하여
$P(T) = \prod_{i = 1, \ldots, d} (T - \alpha_i)$로 쓸 수 있다는
조건과 동치이다. 이 경우 $\alpha$가 $k$ 위에서 {\it 분리 가능}하다고
한다.
\item $\text{d}P/\text{d}T$가 항등적으로 영이다. 이 경우
% P01-E0262
$k$의 표수 $p$는 $ > 0$이고, $P$는 실제로 $T^p$에 관한 다항식이다.
$P$가 $T^q$에 관한 다항식이 되게 하는 가장 큰 거듭제곱
$q = p^e$가 존재함은 분명하다. 그러면 원소 $\alpha^q$는 $k$ 위에서
분리 가능하다.
\end{enumerate}
\end{enumerate}

\begin{definition}
\label{fields-definition-separable-algebraic}
대수적 체확대.
\begin{enumerate}
\item $K$의 모든 원소가 $k$ 위에서 대수적이면 체확대 $K/k$를
{\it 대수적}이라 한다.
\item 모든 $\alpha \in k'$가 $k$ 위에서 분리 가능하면 대수적 확대
$k'/k$를 {\it 분리 가능}이라 한다.
\item $k$의 표수가 $p > 0$이고 모든 원소 $\alpha \in k'$에 대하여
$\alpha^q \in k$가 되게 하는 $p$의 거듭제곱 $q$가 존재하면 대수적
확대 $k'/k$를 {\it 순수 비분리}라 한다.
\item 모든 $\alpha \in k'$에 대하여 $\alpha$의 $k$ 위 최소다항식
$P(T) \in k[T]$가 $k'$ 위에서 일차인수들로 완전히 분해되면 대수적
확대 $k'/k$를 {\it 정규}라 한다.
\item 대수적 확대 $k'/k$가 분리 가능하고 정규이면 {\it 갈루아}라고
한다.
\end{enumerate}
\end{definition}

\noindent
다음 보조정리는 다른 적당한 위치에 들어맞지 않는 듯하다.

\begin{lemma}
\label{fields-lemma-pth-root}
$K$를 표수가 $p > 0$인 체라 하자. $L/K$를 분리 가능 대수적 확대라
하고 $\alpha \in L$이라 하자.
\begin{enumerate}
\item $\alpha$의 $K$ 위 최소다항식의 계수들이 $K$에서 모두
$p$제곱이면 $\alpha$는 $L$에서 $p$제곱이다.
\item 더 일반적으로, $P \in K[T]$가 (a) $\alpha$가 $P$의 근이고,
(b) $P$가 대수적 폐포에서 서로 다른 근들을 가지며, (c) $P$의 모든
계수가 $p$제곱이라는 조건을 만족하는 다항식이면, $\alpha$는 $L$에서
$p$제곱이다.
\end{enumerate}
\end{lemma}

\begin{proof}
정의로부터 (2)가 (1)을 함의한다. $P$가 (2)와 같다고 가정하자.
$P(T) = \sum\nolimits_{i = 0}^d a_i T^{d - i}$ 및 $a_i = b_i^p$로
쓰자. 다항식 $Q(T) = \sum\nolimits_{i = 0}^d b_i T^{d - i}$도
대수적 폐포에서 서로 다른 근들을 갖는다. 실제로 $Q$의 근들은
$P$의 근들의 $p$제곱근이다. $\alpha$가 $p$제곱이 아니면
$T^p - \alpha$는 $L$ 위의 기약다항식이다
(Lemma \ref{fields-lemma-take-pth-root}). 또한 $Q$와 $T^p - \alpha$는
대수적 폐포 $\overline{L}$에서 공통근을 갖는다. 따라서 $Q$와
$T^p - \alpha$는 서로소가 아니며, 이는 $L[T]$에서
$T^p - \alpha | Q$임을 뜻한다. 이는 $Q$의 근들이 서로 다르다는
사실에 모순된다.
\end{proof}
